\documentclass[11pt,oneside,openany]{book}
\usepackage[T1]{fontenc}
\usepackage{lmodern,microtype}
\usepackage[a4paper,margin=28mm]{geometry}
\usepackage{amsmath,amssymb,amsthm,mathtools}
\usepackage[hidelinks,bookmarksopen,bookmarksnumbered]{hyperref}
\hypersetup{pdftitle={Discrete Markov Chains and Martingales},
  pdfauthor={MathGenome},pdfsubject={Unified mathematical notes, P2-07}}
\allowdisplaybreaks[1]
\setlength{\emergencystretch}{2em}

\theoremstyle{plain}
\newtheorem{theorem}{Theorem}[chapter]
\newtheorem{lemma}[theorem]{Lemma}
\newtheorem{proposition}[theorem]{Proposition}
\newtheorem{corollary}[theorem]{Corollary}
\theoremstyle{definition}
\newtheorem{definition}[theorem]{Definition}
\newtheorem{example}[theorem]{Example}
\theoremstyle{remark}
\newtheorem{remark}[theorem]{Remark}

\begin{document}
\frontmatter
\hypersetup{pageanchor=false}
\begin{titlepage}
\thispagestyle{empty}
\vspace*{27mm}
{\large MathGenome \enspace / \enspace P2-07\par}
\vspace{8mm}
\hrule height 0.8pt
\vspace{11mm}
{\Huge\bfseries Discrete Markov Chains\\[3pt] and Martingales\par}
\vspace{4mm}
{\LARGE Recurrence, Mixing, Convergence, and Sampling\par}
\vspace{12mm}
{\large Unified mathematical notes\par}
\vspace{6mm}
{\normalsize Countable chains and classical random walks\\
Discrete martingales, maximal inequalities, and uniform integrability\\
Terminal representations and stopped sums\par}
\vfill
{\large 8 November 2026\par}
\end{titlepage}
\hypersetup{pageanchor=true}
\begingroup
\small
\setlength{\baselineskip}{11.6pt}
\tableofcontents
\endgroup

\chapter*{Scope and conventions}
\addcontentsline{toc}{chapter}{Scope and conventions}

The opening chapter fixes the common language of filtered probability
spaces and stopping times. Chapters 2--4 develop the theory of countable
Markov chains through recurrence, stationarity, periodic and aperiodic
mixing, pathwise time averages, and classical random walks.
Chapters 5--7 develop the martingale theory through Doob inequalities,
uniform integrability, convergence, optional sampling, and Wald
identities. The companion source map connects the forty-three groups
to the twenty-nine original contracts and to their proof dependencies.

Time starts at zero. The terminal \(\sigma\)-algebra of an increasing
filtration is generated by its union. A stopping time may take the
value \(\infty\), and its value there is handled separately in every
statement that evaluates a process at that time. Conditional statements
are almost sure. A Markov chain on a countable set uses the discrete
\(\sigma\)-algebra. We write its transition matrix with the current
state as the row index, and identify probability laws with row vectors
when multiplying by the kernel. Our total variation distance is the
supremum of the absolute difference of the two probabilities over
events. The \(L^p\) norm of a random variable uses the probability
measure in its statement.

The prerequisites are the standard measure and integration interfaces
in \cite{D11}, the construction of path measures, conditional
expectation, and the strong law in \cite{P105}, and the finite-dimensional
real spectral theorem in \cite{D05}. The bounded-variable moment bound
in \cite{P205} is used to connect the independent-variable and
martingale cases. The proofs below make the bridges to this setting
explicit. Textbook references give mathematical provenance and
alternate presentations. Continuous time, general-state Harris
recurrence, scaling limits, and martingale central limit theorems
belong to later units.

\mainmatter
\chapter{Filtrations, Predictability, and Stopping}

\section{Filtrations, adaptation, and martingales}\label{p207:W01}

Work on a probability space \((\Omega,\mathcal F,\mathbb P)\).
We use times \(n\in\mathbb N_0=\{0,1,\ldots\}\). A filtration
\((\mathcal F_n)\) is an increasing sequence of sub-\(\sigma\)-algebras
of \(\mathcal F\). A process \(X=(X_n)\) is adapted if \(X_n\) is
\(\mathcal F_n\)-measurable for every \(n\). Its natural filtration
is \(\sigma(X_0,\ldots,X_n)\). Put
\(\mathcal F_\infty=\sigma(\bigcup_n\mathcal F_n)\).
Write \(\mathbb E[Y\mid\mathcal G]\) for the conditional expectation
of an integrable random variable \(Y\). We use its test-event identity,
linearity, positivity, tower property, and the rule for pulling out
bounded \(\mathcal G\)-measurable factors
\cite[conditional-expectation laws and pull-out rule]{P105}.\footnote{In the
P1-05 LaTeX source the main conditional-expectation identities and
the pull-out rule are labelled \texttt{thm:prob-condexp-laws}
and \texttt{thm:prob-condexp-pullout}, respectively. The companion
map gives exact source locations.}
Conditional equalities and inequalities hold almost surely. Countably
many such relations can be used together on a single full-probability
event by removing the union of their exceptional sets.

\begin{definition}
A real adapted process with \(\mathbb E|X_n|<\infty\) for every \(n\) is a
\emph{martingale} if
\(\mathbb E[X_{n+1}\mid\mathcal F_n]=X_n\) for every \(n\).
It is a \emph{submartingale} if
\(\mathbb E[X_{n+1}\mid\mathcal F_n]\ge X_n\) for every \(n\).
\end{definition}

\begin{proposition}
The one-step condition is equivalent to the same equality, or
inequality, with \(X_{n+1},\mathcal F_n,X_n\) replaced by
\(X_j,\mathcal F_i,X_i\) for every \(i\le j\).
\end{proposition}
\begin{proof}
Fix \(i\) and induct on \(j\ge i\). For \(j=i\), adaptation and
integrability give the conditional identity. For \(j+1\), take
\(\mathbb E[\cdot\mid\mathcal F_i]\) in the one-step relation at
\(j\). By the tower property its left side is
\(\mathbb E[X_{j+1}\mid\mathcal F_i]\). Its right side is
\(\mathbb E[X_j\mid\mathcal F_i]\), so the induction hypothesis
applies. The converse uses \(j=i+1\).
\end{proof}

\begin{example}
If \(\xi_k\) is integrable and \(\mathcal F_k\)-measurable,
\(\mathbb E\xi_k=0\), and \(\xi_k\) is independent of
\(\mathcal F_{k-1}\) for \(k\ge1\), then
\(M_n=\sum_{k=1}^n\xi_k\) is a martingale. Independence gives
\(\mathbb E[\xi_k1_A]=0\) for \(A\in\mathcal F_{k-1}\), so the
conditional expectation of the increment is zero. Also, any
\(Y\in L^1\) defines the martingale \(M_n=\mathbb E[Y\mid\mathcal F_n]\).
This follows from adaptation and the tower property.
\end{example}

The terminology follows \cite[Sections 10.1--10.6]{Williams}
and \cite[Chapter 9]{Kallenberg}.

\section{Stopping times and the past at a stopping time}\label{p207:W02}

\begin{definition}
A stopping time is a map
\(\tau:\Omega\to\mathbb N_0\cup\{\infty\}\) such that
\(\{\tau\le n\}\in\mathcal F_n\) for every finite \(n\).
Set
\[
 \mathcal F_\tau=\{A\in\mathcal F_\infty:
                 A\cap\{\tau\le n\}\in\mathcal F_n
                 \text{ for every finite }n\}.
\]
The stopped process is \(X^\tau_n=X_{n\wedge\tau}\).
The stopped value \(X_\tau\) means \(X_n\) on \(\{\tau=n\}\).
On \(\{\tau=\infty\}\) its value is specified separately whenever a
full-space random variable is needed.
\end{definition}

\begin{proposition}
The collection \(\mathcal F_\tau\) is a \(\sigma\)-algebra.
For \(A\in\mathcal F_\infty\), membership in it is equivalent
to \(A\cap\{\tau=n\}\in\mathcal F_n\) for every finite \(n\).
If \(\sigma\le\tau\) pointwise, then
\(\mathcal F_\sigma\subseteq\mathcal F_\tau\).
If \(X\) is adapted, then \(X^\tau\) is adapted. If each \(X_n\)
is integrable, so is each \(X^\tau_n\). For adapted \(X\) and an
\(\mathcal F_\infty\)-measurable random variable \(Y\), the version
of \(X_\tau\) that equals \(Y\) on \(\{\tau=\infty\}\) is
\(\mathcal F_\tau\)-measurable.
\end{proposition}
\begin{proof}
Closure under complements follows from
\(A^c\cap\{\tau\le n\}=\{\tau\le n\}\setminus
(A\cap\{\tau\le n\})\), and countable unions distribute over the
intersection. Difference of consecutive cutoff events gives
\(\{\tau=n\}\in\mathcal F_n\), with \(\{\tau=0\}=\{\tau\le0\}\).
Thus \(A\cap\{\tau=n\}\) is the difference of the corresponding
two cutoffs. Conversely, the finite union of layers through \(n\)
belongs to \(\mathcal F_n\). If \(\sigma\le\tau\) and
\(A\in\mathcal F_\sigma\), then
\[
 A\cap\{\tau\le n\}
   =(A\cap\{\sigma\le n\})\cap\{\tau\le n\}\in\mathcal F_n .
\]

At time \(n\), partition \(\Omega\) into the events \(\{\tau=k\}\)
for \(k<n\) and \(\{\tau\ge n\}\). All are in \(\mathcal F_n\),
and the stopped process equals the corresponding \(X_k\) or \(X_n\).
It is therefore measurable there and bounded in absolute value by
\(\sum_{k=0}^n |X_k|\). For the full-space stopped value use the
countable partition into finite layers and \(\{\tau=\infty\}\), all
in \(\mathcal F_\infty\). The inverse image of a Borel set on each
finite layer is in \(\mathcal F_n\), so the layer characterization
proves the last claim.
\end{proof}

In particular \(\{\tau<\infty\}\in\mathcal F_\tau\).
Also \(\mathcal F_{n}=\mathcal F_\tau\) for constant \(\tau=n\):
the defining condition at \(n\) gives inclusion in \(\mathcal F_n\),
and every event in \(\mathcal F_n\) satisfies all cutoffs.
See \cite[Section 10.9 and Appendix A14]{Williams}.

\section{Predictable transforms and the Doob decomposition}\label{p207:W03}

Write \(\Delta X_k=X_{k+1}-X_k\).
The coefficient \(H_k\) used on this increment is predictable when
it is \(\mathcal F_k\)-measurable. The time-\(n\) gain from a predictable
sequence \(H\) is
\[
 (H\mathbin{\cdot}X)_n=\sum_{k=0}^{n-1}H_k\Delta X_k .
\]

\begin{proposition}
If every \(H_k\) is bounded, this transform is a martingale for
martingale \(X\). If \(X\) is a submartingale and \(H_k\ge0\), it is
a submartingale. Both conclusions also hold when the individual
products \(H_k\Delta X_k\) are integrable, with finite real \(H_k\).
\end{proposition}
\begin{proof}
Each term in the finite sum at time \(n\) is integrable and
\(\mathcal F_n\)-measurable. For bounded \(H_k\), the conditional
expectation of its increment is
\(H_k\mathbb E[\Delta X_k\mid\mathcal F_k]\). It is zero in the
martingale case and nonnegative in the submartingale case.
For an unbounded coefficient, replace it by
\(H_k1_{\{|H_k|\le m\}}\). The products converge in \(L^1\) by
domination by \(|H_k\Delta X_k|\). Thus their conditional
expectations converge in \(L^1\). Zero and nonnegativity are
preserved, the latter because the negative part of the limit has
\(L^1\) norm at most the \(L^1\) distance to each nonnegative
approximant.
\end{proof}

\begin{theorem}[Discrete Doob decomposition]
Every integrable real adapted process \(X\) has a decomposition
\(X_n=M_n+A_n\), where \(M\) is a martingale, \(A_0=0\), \(A\)
is integrable, and \(A_{n+1}\) is \(\mathcal F_n\)-measurable.
It is unique up to almost-sure equality at every time. Moreover
\(X\) is a submartingale if and only if \(A\) is nondecreasing
almost surely.
\end{theorem}
\begin{proof}
Define \(A_n=\sum_{k<n}\mathbb E[\Delta X_k\mid\mathcal F_k]\)
using measurable representatives. The summands are integrable
and the sum at \(n+1\) is \(\mathcal F_n\)-measurable. Set
\(M_n=X_n-A_n\). Its conditional expected increment is zero
by construction. Any decomposition with the indicated
predictability must satisfy
\(A_{n+1}-A_n=\mathbb E[\Delta X_n\mid\mathcal F_n]\) almost
surely because the martingale increment contributes zero.
Thus the increments and the common initial \(A_0=0\) determine
the process. Finally these increments are nonnegative precisely
when \(X\) is a submartingale.
\end{proof}

Compare \cite[Sections 10.6--10.8 and 12.11]{Williams} and
\cite[Lemma 9.10]{Kallenberg}.

\section{Stopping preserves martingales and submartingales}\label{p207:W04}

\begin{theorem}
For any stopping time \(\tau\), \(X^\tau\) is a martingale
when \(X\) is a martingale and a submartingale when \(X\) is a
submartingale.
\end{theorem}
\begin{proof}
The stopped process is integrable and adapted by
Section~\ref{p207:W02}. On \(\{\tau\le k\}\) its increment at
\(k\) is zero. On \(\{\tau>k\}\) it is \(\Delta X_k\).
Thus the stopped increments form the predictable transform
with \(H_k=1_{\{\tau>k\}}\), which is bounded and nonnegative.
Apply Section~\ref{p207:W03} and add the common integrable
\(\mathcal F_0\)-measurable initial value \(X_0\).
\end{proof}
\chapter{Countable Markov Chains and Recurrence}

\section{Countable chains and transition powers}\label{p207:W05}

Let \(S\) be a nonempty countable set with its full sigma-algebra.
For \(i,j\in S\), write \(\delta_{ij}=\mathbf{1}_{\{i=j\}}\).
A stochastic matrix \(P=(P_{ij})_{i,j\in S}\) has nonnegative entries
and \(\sum_j P_{ij}=1\) for every \(i\). For \(A\subseteq S\) and a
nonnegative function \(f\) on \(S\), put
\[
 P(i,A)=\sum_{j\in A}P_{ij},
 \qquad
 (Pf)(i)=\sum_j P_{ij}f(j).
\]
In extended nonnegative sums we use \(0\cdot\infty=0\).
Products of stochastic matrices are defined by
\((PQ)_{ij}=\sum_xP_{ix}Q_{xj}\). Nonnegative summation shows that a
product is stochastic and that this multiplication is associative.
Let \(P^0=I=(\delta_{ij})\) and \(P^{m+1}=P^mP\). If \(\lambda\) is
a probability on \(S\), the probability \(\lambda P^m\) has masses
\((\lambda P^m)_j=\sum_i\lambda_i(P^m)_{ij}\).

\begin{definition}
A time-homogeneous \(P\)-Markov chain with initial law \(\lambda\)
relative to \((\mathcal F_n)\) is an \(S\)-valued adapted process
\((X_n)_{n\ge0}\) with \(\mathbb P(X_0=i)=\lambda_i\) and
\[
 \mathbb E[\mathbf{1}_{\{X_{n+1}=j\}}\mid\mathcal F_n]=P_{X_nj}
 \quad\text{almost surely}
\]
for every \(n\ge0\) and \(j\in S\).
\end{definition}

The natural filtration is \(\sigma(X_0,\ldots,X_n)\) at time \(n\).
With a larger filtration the one-step property is an assumption
on that filtration. All conditional identities below refer to the
filtration with respect to which the chain is Markov.

\begin{theorem}[Transition probabilities at deterministic times]
For \(m,n\ge0\) and \(A\subseteq S\),
\[
 \mathbb E[\mathbf{1}_{\{X_{n+m}\in A\}}\mid\mathcal F_n]
     =P^m(X_n,A)\quad\text{almost surely}.
\]
In particular \(X_n\) has law \(\lambda P^n\).
\end{theorem}
\begin{proof}
For a singleton \(A=\{j\}\), the case \(m=0\) follows from
adaptedness. Suppose the result holds at \(m\). The tower property
and the one-step identity express the conditional expectation at
\(m+1\) as \(\mathbb E[P_{X_{n+m}j}\mid\mathcal F_n]\).
Expand the random variable into the countable sum
\(\sum_x P_{xj}\mathbf{1}_{\{X_{n+m}=x\}}\).
Conditional monotone convergence and the induction hypothesis
give \(\sum_x(P^m)_{X_nx}P_{xj}=(P^{m+1})_{X_nj}\).
Sum over \(j\in A\) by the same monotone convergence argument.
At \(n=0\), taking the expectation of the \(m\)-step formula
gives the marginal law of \(X_m\).
\end{proof}

These conventions agree with \cite[Sections 1.1--1.2]{Norris}
and \cite[Section 5.2]{Durrett}.

\section{Path laws and Chapman--Kolmogorov}\label{p207:W06}

Let \(\mathbb P_\lambda\) denote the law on \(S^{\mathbb N_0}\) of a
chain with initial law \(\lambda\) and transition matrix \(P\), as
constructed below. We write \(\mathbb P_i\) and \(\mathbb E_i\) for
probability and expectation when \(\lambda\) is the point mass at \(i\).
The path space always carries its product sigma-algebra.

\begin{theorem}[Canonical law]
There is a unique probability on \(S^{\mathbb N_0}\) for which every
atomic prefix cylinder has probability
\[
 \mathbb P_\lambda(X_0=i_0,\ldots,X_n=i_n)
    =\lambda_{i_0}\prod_{k=0}^{n-1}P_{i_k i_{k+1}}.
 \tag{*}
\]
The coordinate process is \(P\)-Markov relative to the natural
filtration. Every \(P\)-Markov chain with initial law \(\lambda\) has
this joint path law. The transition matrices satisfy
\[
 (P^{m+n})_{ij}=\sum_x(P^m)_{ix}(P^n)_{xj}
 \qquad (m,n\ge0).
\]
\end{theorem}
\begin{proof}
For a history \(x_0,\ldots,x_n\), the probability kernel for the next
coordinate is \(A\mapsto P(x_n,A)\). It is measurable in the history
because the finite history space is countable and discrete.
Apply the Ionescu--Tulcea extension theorem
\cite[Theorem \texttt{thm:ionescu-tulcea}]{P105} with initial law
\(\lambda\) and these history kernels.
Integrating the indicator of an atomic prefix against the successive
kernels multiplies the previous prefix probability by \(P_{i_n i_{n+1}}\).
This proves \((*)\) by induction, with empty product one at \(n=0\).

At fixed \(n\), the prefix atoms form a countable partition. On each
atom \(B=\{X_0=i_0,\ldots,X_n=i_n\}\), equation \((*)\) gives
\(\mathbb P_\lambda(B\cap\{X_{n+1}=j\})
 =P_{i_nj}\mathbb P_\lambda(B)\). Adding these identities over the
atoms in a natural-filtration event proves the conditional identity.
Conversely, in any \(P\)-Markov chain, the same identity with
\(B\in\mathcal F_n\) proves \((*)\) recursively. Two laws satisfying
\((*)\) thus agree on all finite prefix cylinders, by adding over
prefix atoms. These cylinders form a generating \(\pi\)-system, so
probability uniqueness gives equality on the product sigma-algebra.

For the last formula, associativity proved in
Section~\ref{p207:W05} gives \(P^{m+n}=P^mP^n\) by induction on \(n\),
including \(P^0=I\). Taking the \((i,j)\) entry proves the identity.
\end{proof}

For the classical countable-chain construction see
\cite[Section 1.1]{Norris}.

\section{The strong Markov property}\label{p207:W07}

For bounded measurable \(g:S^{\mathbb N_0}\to\mathbb R\), let
\(G(i)=\mathbb E_i[g(X_0,X_1,\ldots)]\). The function \(G\) is bounded
and measurable on \(S\).

\begin{theorem}[Conditional future law]
Let \(\tau\) be an \(\mathbb N_0\cup\{\infty\}\)-valued stopping time of a
\(P\)-Markov chain. For each \(A\in\mathcal F_\tau\),
\[
 \int_{A\cap\{\tau<\infty\}}g(X_\tau,X_{\tau+1},\ldots)\,d\mathbb P
 =
 \int_{A\cap\{\tau<\infty\}}G(X_\tau)\,d\mathbb P.
 \tag{**}
\]
In particular, when \(\tau\) is finite almost surely, the conditional
law given \(\mathcal F_\tau\) of the shifted path is the canonical law
started from \(X_\tau\).
\end{theorem}
\begin{proof}
First fix \(n\) and an atomic prefix cylinder of the future
\(X_n,\ldots,X_{n+m}\). Conditioning first at time \(n+m-1\),
the indicator of \(\{X_n=i_0,\ldots,X_{n+m-1}=i_{m-1}\}\)
is measurable there and the conditional factor for the last
coordinate is \(P_{X_{n+m-1},i_m}\). On that prefix event the
factor is \(P_{i_{m-1},i_m}\).
Iterating gives the conditional probability given \(\mathcal F_n\)
as \(\mathbf{1}_{\{X_n=i_0\}}\prod_{k<m}P_{i_ki_{k+1}}\),
exactly the value of \(G(X_n)\) for this indicator. This also
holds for an arbitrary prefix cylinder by summation over atoms.
For each fixed \(B\in\mathcal F_n\), the integral of the shifted
indicator over \(B\) and the integral over \(B\) of its canonical
probability given \(X_n\) define two finite measures on the path
space with the same total mass. They agree on prefix cylinders
and hence on all measurable sets. Approximation by simple
functions proves the fixed-time integral identity for bounded \(g\).

For \(A\in\mathcal F_\tau\), put \(A_n=A\cap\{\tau=n\}\in\mathcal F_n\).
Apply the fixed-time identity to \(A_n\), then sum over \(n\).
The integrals are absolutely summable because the \(A_n\) are
disjoint and \(g,G\) are bounded. This is \((**)\).
On the finite stopping event, \(G(X_\tau)\) is stopped-measurable:
the intersection of a level set with \(\{\tau\le n\}\) is the
union of its time-\(k\) pieces for \(k\le n\).
It is \(\mathcal F_\infty\)-measurable after assigning a fixed value
on \(\{\tau=\infty\}\), which suffices for membership in
\(\mathcal F_\tau\). The same countable decomposition gives
measurability of the future path on \(\{\tau<\infty\}\).
Thus \((**)\) characterizes the stated conditional law, with the
finite-event restriction when needed.
\end{proof}

For the finite-stopping-time formulation see
\cite[Theorem 1.4.2 and Section 6.5]{Norris}.
For nonnegative measurable \(g\), \((**)\) extends by truncation and
monotone convergence, with integrals in the extended nonnegative reals.

\section{First-hit equations and minimal solutions}\label{p207:W08}

Let \(T_A=\inf\{n\ge0:X_n\in A\}\), with infimum \(\infty\) when
the set is empty. Define
\(h(i)=\mathbb P_i(T_A<\infty)\) and \(u(i)=\mathbb E_iT_A\).

\begin{theorem}[Minimal solutions]
The function \(h\) is the smallest nonnegative solution of
\[
 h=1\ \text{on }A,\qquad h=Ph\ \text{on }S\setminus A.
\]
The function \(u\) is the smallest \([0,\infty]\)-valued solution of
\[
 u=0\ \text{on }A,\qquad u=1+Pu\ \text{on }S\setminus A.
\]
\end{theorem}
\begin{proof}
Put \(h_N(i)=\mathbb P_i(T_A\le N)\) and
\(u_N(i)=\mathbb E_i(T_A\wedge N)\). At \(N=0\) these are
\(\mathbf{1}_A\) and \(0\). On \(A\), they retain these values
for all \(N\). Outside \(A\), first-step conditioning gives
\(h_{N+1}=Ph_N\) and \(u_{N+1}=1+Pu_N\). In the second equality,
the identity for the duration is
\((1+T_A\text{ of the shifted path})\wedge(N+1)
  =1+(T_A\text{ of the shifted path}\wedge N)\).
Both sequences are increasing. Taking their limits and using
monotone convergence under the first-step kernel proves the
equations. Any nonnegative solution dominates the \(N=0\)
function and, by the equations, dominates every later iterate.
It therefore dominates the limit.
\end{proof}

This argument works also for \(A=\varnothing\), when \(h=0\) and
\(u=\infty\). Compare \cite[Theorems 1.3.2 and 1.3.5]{Norris}.

\section{Communication, return, and recurrence}\label{p207:W09}

State \(j\) is accessible from \(i\) if \((P^n)_{ij}>0\) for some
\(n\ge0\). Two states communicate if each is accessible from the
other. This is an equivalence relation, since positive transitions
along finitely many steps compose by Chapman--Kolmogorov. The chain
is irreducible when \(S\) is a single communicating class.
A class is closed if no state in it has a positive one-step
transition to its complement.

The first positive return to \(i\) is
\(\tau_i^+=\inf\{n\ge1:X_n=i\}\). This is a stopping time, as
\(\{\tau_i^+\le n\}\) is a finite union of state events at times
\(1,\ldots,n\). With \(r_i=\mathbb P_i(\tau_i^+<\infty)\), state \(i\)
is recurrent if \(r_i=1\) and transient if \(r_i<1\).
It is positive recurrent if \(\mathbb E_i\tau_i^+<\infty\).
Finite extended expectation implies recurrence, since any positive
mass at infinity makes the expectation infinite.
A recurrent state with infinite mean return time is null recurrent.
These definitions follow \cite[Sections 1.2, 1.5 and 1.7]{Norris}.

\section{The Green series and number of visits}\label{p207:W10}

\begin{theorem}[Green criterion]
For a chain started at \(i\), let \(N_i\) be the number of visits
to \(i\), including time zero. For every integer \(k\ge1\),
\[
 N_i=\sum_{n\ge0}\mathbf{1}_{\{X_n=i\}},\qquad
 \mathbb P_i(N_i\ge k)=r_i^{\,k-1},
\]
where \(r_i^0=1\).
In particular,
\[
 \sum_{n\ge0}(P^n)_{ii}=\mathbb E_iN_i
 =\begin{cases}
 (1-r_i)^{-1},& r_i<1,\\
 \infty,& r_i=1.
 \end{cases}
\]
A recurrent state is visited infinitely often almost surely when the
chain starts there.
\end{theorem}
\begin{proof}
Let \(R_0=0\) and \(R_{k+1}=\inf\{n>R_k:X_n=i\}\), with infinity
absorbing. These are stopping times: for \(R_{k+1}=n<\infty\),
partition by \(R_k=m<n\) and require no visit at times \(m+1,\ldots,n-1\)
and a visit at \(n\). Each resulting event lies in \(\mathcal F_n\).
On \(\{R_k<\infty\}\), the conditional probability of another visit
is \(r_i\) by the strong Markov property. Thus
\(\mathbb P_i(R_k<\infty)=r_i^k\) by induction.
The event \(N_i\ge k\) equals \(\{R_{k-1}<\infty\}\).
Summing the tail probabilities gives the displayed expectation
and \(\mathbb P_i(N_i=\infty)=\lim_k r_i^{k-1}\).
The equality with the transition sum follows from Tonelli applied
to the indicators defining \(N_i\).
\end{proof}

See \cite[Section 1.5]{Norris} and \cite[Section 5.3]{Durrett}.

\section{First-return convolution and generating functions}\label{p207:W11}

\begin{proposition}
Fix \(i\), set \(p_n=(P^n)_{ii}\) for \(n\ge0\) and
\(f_n=\mathbb P_i(\tau_i^+=n)\) for \(n\ge1\). Then \(p_0=1\) and
\(p_n=\sum_{k=1}^n f_kp_{n-k}\) for \(n\ge1\).
For \(0\le z<1\), the generating functions
\(G(z)=\sum_{n\ge0}p_nz^n\) and \(F(z)=\sum_{n\ge1}f_nz^n\) satisfy
\[
 G(z)=\frac{1}{1-F(z)},\qquad
 \lim_{z\uparrow1}F(z)=r_i,\qquad
 \lim_{z\uparrow1}G(z)=\sum_{n\ge0}p_n.
\]
\end{proposition}
\begin{proof}
On \(\{X_n=i\}\) for \(n\ge1\), the first positive return has some
value \(k\le n\). By the strong Markov property on
\(\{\tau_i^+=k\}\), the probability of also being at \(i\) at time
\(n\) is \(f_kp_{n-k}\). Summing yields the recursion.
Multiplying by \(z^n\), summing and using nonnegative summation
gives \(G(z)=1+F(z)G(z)\). Here \(G(z)\le(1-z)^{-1}<\infty\)
and \(F(z)\le z\sum_{n\ge1}f_n\le z<1\).
The limits follow by monotone convergence. The resulting value
of \(G(1-)\) agrees with the Green criterion in
Section~\ref{p207:W10}, including divergence if \(r_i=1\).
\end{proof}

The first-return decomposition is the one used in
\cite[Section 1.5]{Norris} and \cite[Section 5.3]{Durrett}.

\section{Classes, periods, and recurrence}
\label{p207:W12}

Write \(d(i)=\gcd\{n\ge1:(P^n)_{ii}>0\}\), with \(d(i)=0\) if this
set is empty. A state or class of period one is called aperiodic.

\begin{theorem}
Communicating states have the same recurrence classification
(recurrent or transient) and the same period. A communicating class
containing a recurrent state is closed. Within a recurrent
communicating class the first hit of any specified state is finite
almost surely, from every start in the class. Every finite
irreducible chain is positive recurrent.
\end{theorem}
\begin{proof}
If \(i,j\) communicate, choose lengths \(a,b\ge0\) with positive
transition probabilities from \(i\) to \(j\) and from \(j\) to \(i\),
respectively. Chapman--Kolmogorov gives
\[
 (P^{a+n+b})_{ii}\ \ge\ (P^a)_{ij}(P^n)_{jj}(P^b)_{ji}.
\]
By the Green criterion, recurrence at \(j\) implies recurrence at
\(i\). The symmetric inequality gives the converse.
For distinct communicating states, \(a+b>0\), and this length is
a return length at both states. If \(n\) is a return length at \(j\),
the displayed inequality makes \(a+n+b\) a return length at \(i\).
Their difference shows that \(d(i)\) divides \(n\). Thus \(d(i)\)
divides \(d(j)\) and conversely. For equal states the assertion
holds by definition. This also covers the empty return-set
convention, since distinct communicating states have a common
positive return length.

Positive probability of a finite transition is equivalent to the
existence of a path of that length whose step probabilities are all
positive, by expanding the finite matrix power as a countable sum
of nonnegative products. If \(i\) is recurrent and \(j\ne i\) is
accessible from \(i\), take a shortest positive-probability path
from \(i\) to \(j\). It has no intermediate visit to \(i\).
If \(i\) were inaccessible from \(j\), a chain started at \(j\)
would have probability zero of ever visiting \(i\), by summing the
probabilities at deterministic times. Following the chosen path
from \(i\) would then give positive probability of no return to \(i\),
a contradiction. Thus all states accessible from \(i\) communicate
with \(i\). Any state reached by one step from a state in its class
is accessible from \(i\), proving closedness.

If \(i\ne j\) lie in that class, the same shortest-path argument
gives \(\theta=\mathbb P_i(T_{\{j\}}<\tau_i^+)>0\).
All successive return times to \(i\) are finite almost surely.
At each return, the strong Markov property gives probability
\(\theta\) of a visit to \(j\) before the next return, conditional
on the stopped past. The event of having missed \(j\) in earlier
excursions is measurable in that past, as is seen by splitting
at the finite return times. Inductively the probability of missing
\(j\) through \(k\) excursions is \((1-\theta)^k\), which tends
to zero. For \(i=j\) the first hit includes time zero.

Finally suppose \(S\) is finite and the chain is irreducible.
Fix \(j\). For every starting state choose a positive-probability
path to \(j\), of length zero if already at \(j\). Let \(K\ge1\)
bound the lengths and let \(\epsilon>0\) be the minimum of their
path probabilities. Then \(\mathbb P_i(T_{\{j\}}\le K)\ge\epsilon\)
uniformly in \(i\). Conditioning at successive deterministic times
\(kK\), the probability of having missed \(j\) by \(kK\) is at
most \((1-\epsilon)^k\). The expectation of \(T_{\{j\}}\) is
therefore bounded by \(K/\epsilon\). Starting at \(j\) and taking
one step first gives \(\mathbb E_j\tau_j^+\le 1+K/\epsilon\).
Since \(j\) was arbitrary, every state is positive recurrent.
\end{proof}

See \cite[Sections 1.2, 1.5 and 1.7--1.8]{Norris}.
\chapter[Stationarity and Mixing]{Stationarity, Mixing, and Long-Run Averages}

\section{Stationarity, detailed balance, and reversal}\label{p207:W13}

A probability \(\pi\) on \(S\) is stationary for \(P\) if
\(\pi P=\pi\). It satisfies detailed balance if
\(\pi_i P_{ij}=\pi_jP_{ji}\) for all \(i,j\).
A stationary chain is reversible if each finite path segment has
the same law as the segment read in the opposite order.

\begin{proposition}
A chain with initial law \(\pi\) has shift-invariant finite-dimensional
distributions if and only if \(\pi P=\pi\).
Detailed balance implies stationarity.
If \(P\) is irreducible and \(\pi\) is stationary, then \(\pi_i>0\)
for all \(i\).
For any stationary \(\pi\), the reversed segment has the law of a
chain started from \(\pi\) with kernel
\[
 P^*_{ij}=\frac{\pi_jP_{ji}}{\pi_i}
\]
on the support \(\{i:\pi_i>0\}\). Reversibility is equivalent to
detailed balance.
\end{proposition}
\begin{proof}
If \(\pi P=\pi\), all marginal laws are \(\pi\).
The deterministic-time Markov property and the path product
in Section~\ref{p207:W06} give the same consecutive-block law
at every starting time, hence also the same laws of arbitrary
finite subsequences after a shift. The converse follows by comparing
the marginals at times zero and one. Summing detailed balance over
\(i\) gives \((\pi P)_j=\pi_j\sum_iP_{ji}=\pi_j\).
For irreducible \(P\), choose \(i\) with \(\pi_i>0\).
For each \(j\), there is an \(m\) with \((P^m)_{ij}>0\).
Stationarity gives \(\pi_j\ge\pi_i(P^m)_{ij}>0\).

In general the support is closed, since for \(j\) off the support
\(\sum_i\pi_iP_{ij}=0\) forces \(P_{ij}=0\) for each supported \(i\).
On the support, stationarity makes the rows of \(P^*\) sum to one.
For states \(i_0,\ldots,i_n\) on the support,
\[
 \pi_{i_0}\prod_{k<n}P^*_{i_k i_{k+1}}
 =\pi_{i_n}\prod_{k<n}P_{i_{k+1}i_k}.
\]
The right side is the forward probability of the reversed word.
This identifies all reversed segment laws. Detailed balance makes
\(P^*=P\) on the support, hence gives reversibility. Conversely,
reversibility on segments of length one step equates the two
probabilities \(\pi_iP_{ij}\) and \(\pi_jP_{ji}\) for every pair.
\end{proof}

Compare \cite[Sections 1.7 and 1.9]{Norris}
and \cite[Section 5.5]{Durrett}.

\section{Stationarity and mean return times}
\label{p207:W14}

\begin{theorem}[Stationary probability on an irreducible class]
A countable irreducible chain has a stationary probability
if and only if some state is positive recurrent. In that case
all states are positive recurrent, the probability \(\pi\) is
unique, and
\[
 \pi_i=\frac{1}{\mathbb E_i\tau_i^+}
 \qquad(i\in S).
\]
For each \(i\), the expected occupation in an excursion from \(i\)
back to \(i\) is \(\pi_j/\pi_i\) at state \(j\).
\end{theorem}
\begin{proof}
First assume \(\pi\) is stationary. For fixed \(i\), the events
\[
 L_n=\{X_n=i,\ X_k\ne i\text{ for all }k>n\},
 \qquad n\ge0,
\]
are pairwise disjoint under \(\mathbb P_\pi\).
By the fixed-time conditional future law in Section~\ref{p207:W07},
each has probability \(\pi_i(1-r_i)\).
Their sum is at most one, and \(\pi_i>0\) by irreducibility,
so \(r_i=1\).

Now for a recurrent \(i\), define its excursion occupation masses
\[
 \gamma_j=\mathbb E_i\left[\sum_{0\le n<\tau_i^+}
                             \mathbf{1}_{\{X_n=j\}}\right].
\]
We have \(\gamma_i=1\). For \(j\ne i\), a shortest path from \(j\)
to \(i\) has no intermediate visit to \(j\). Its positive probability
gives \(\theta=\mathbb P_j(T_{\{i\}}<\tau_j^+)>0\).
Whenever the excursion reaches its \(k\)th visit to \(j\), the
conditional probability of reaching the next one before hitting \(i\)
is at most \(1-\theta\) by the strong Markov property.
Indeed the visit time is a stopping time by the same decomposition
as for successive returns, and the event that it precedes
\(\tau_i^+\) is measurable there, as can be checked on each finite
visit-time layer. Thus the probability of at least \(k\) visits
to \(j\) in the excursion is at most \((1-\theta)^{k-1}\).
Its tail sum shows \(\gamma_j<\infty\).

Nonnegative summation and conditioning at time \(n\) yield
\[
 \begin{aligned}
 \sum_{k}\gamma_kP_{kj}
 &=\sum_{n\ge0}\mathbb E_i[
       \mathbf{1}_{\{n<\tau_i^+\}}P_{X_nj}]\\
 &=\mathbb E_i\!\left[\sum_{0\le n<\tau_i^+}
                         \mathbf{1}_{\{X_{n+1}=j\}}\right]
 =\gamma_j .
 \end{aligned}
\]
In the middle equality the stopping event belongs to
\(\mathcal F_n\). The last equality uses \(X_0=X_{\tau_i^+}=i\)
and the almost-sure finiteness of \(\tau_i^+\).
Moreover \(\sum_j\gamma_j=\mathbb E_i\tau_i^+\), by Tonelli.
If this total mass is finite, normalization gives a stationary
probability.

To prove the converse and uniqueness, let \(\pi\) be any stationary
probability and put \(\nu=\pi/\pi_i\).
For \(j\ne i\), let \(a_j=P_{ij}\) and let \(Q=(P_{kj})_{k,j\ne i}\)
be the matrix of transitions avoiding \(i\).
Writing \(\nu'=(\nu_j)_{j\ne i}\), stationarity and \(\nu_i=1\) give
\(\nu'=a+\nu' Q\). Iterating this equation with nonnegative terms
gives \(\nu_j\ge \sum_{m=0}^{M-1}(aQ^m)_j\) for every \(M\ge1\).
The term \((aQ^m)_j\) is the probability of being at \(j\) at
time \(m+1\) with no intervening return to \(i\), by the path
formula. Taking \(M\to\infty\) shows \(\nu_j\ge\gamma_j\).
Equality at \(i\) holds by normalization. Consequently the
coordinatewise difference \(\eta=\nu-\gamma\) is a nonnegative
invariant measure with \(\eta_i=0\). Subtraction in the invariance
equations is legitimate since both sums in each coordinate are
finite. By iteration and Tonelli,
\[
 0=\eta_i=\sum_k\eta_k(P^m)_{ki}
 \quad (m\ge0).
\]
Irreducibility gives for each \(k\) an \(m\) with a positive factor,
so \(\eta=0\). Therefore \(\pi_i^{-1}=\sum_j\gamma_j
=\mathbb E_i\tau_i^+\), and \(\pi\) is the unique normalization
of the excursion measure. Since the argument works for every
reference state, it proves positive recurrence and the occupation
formula for all states. It also proves the implication from one
positive recurrent state by using the normalization constructed above.
\end{proof}

A closed communicating class can be treated by restricting \(P\)
to that class. In particular positive and null recurrence are
class properties within recurrent classes. See
\cite[Theorems 1.7.5--1.7.7]{Norris} and
\cite[Section 5.5]{Durrett}.

\section{Total variation and coupling}\label{p207:W15}

\begin{proposition}
For probabilities \(\mu,\nu\) on a countable set, define their
total variation distance by
\[
 d_{\mathrm{TV}}(\mu,\nu)=\sup_{A\subseteq S}|\mu(A)-\nu(A)|.
\]
Then
\[
 d_{\mathrm{TV}}(\mu,\nu)
   =\frac12\sum_i|\mu_i-\nu_i|
   =1-\sum_i\min(\mu_i,\nu_i).
\]
A stochastic matrix \(P\) is a contraction for this distance.
If \((U,V)\) has marginals \(\mu,\nu\), then
\(d_{\mathrm{TV}}(\mu,\nu)\le \mathbb P(U\ne V)\), and equality is
attained by some coupling. If \(\mu_n,\mu\) are probability laws
and \(\mu_n(i)\to\mu_i\) for every \(i\), then
\(d_{\mathrm{TV}}(\mu_n,\mu)\to0\).
\end{proposition}
\begin{proof}
The positive and negative parts of the signed masses \(\mu_i-\nu_i\)
have equal sum because the total difference is zero. The supremum
is that common sum, attained by taking the set where
\(\mu_i\ge\nu_i\). This gives the first formula, and
\(\min(\mu_i,\nu_i)=(\mu_i+\nu_i-|\mu_i-\nu_i|)/2\) gives the second.
For contraction, sum the pointwise bound
\[
 |(\mu P)_j-(\nu P)_j|
 \le\sum_i|\mu_i-\nu_i|P_{ij}
\]
over \(j\), then use the row sums.
For any coupling, the difference of the two indicators of \(A\)
vanishes on \(U=V\) and has absolute value at most one.
Taking expectations and the supremum proves the coupling bound.
To attain it, put mass \(\min(\mu_i,\nu_i)\) on each \((i,i)\).
If the remaining mass is positive, use the product of the
normalized remaining marginals with that remaining weight.
Their supports are disjoint, so this part has \(U\ne V\)
almost surely. If no mass remains, the diagonal suffices.
For the last statement, choose a finite set of \(\mu\)-mass
greater than \(1-\epsilon\). On that set the sum of common masses
tends to its \(\mu\)-mass. The common-mass formula gives a
\(\limsup\) distance of at most \(\epsilon\), which is arbitrary.
\end{proof}

See \cite[Section 1.8]{Norris} and \cite[Section 5.6]{Durrett}.

\section{Aperiodic convergence in total variation}\label{p207:W16}

\begin{theorem}
If \(P\) is irreducible, positive recurrent and aperiodic, its
stationary probability \(\pi\) satisfies
\[
 d_{\mathrm{TV}}\big(P^n(i,\cdot),\pi\big)\longrightarrow 0
 \quad (i\in S).
\]
For every initial probability \(\lambda\), the same conclusion
holds with \(\lambda P^n\) in place of \(P^n(i,\cdot)\).
\end{theorem}
\begin{proof}
Fix \(a\in S\). The positive return lengths at \(a\) are closed
under addition by Chapman--Kolmogorov. Their gcd is one.
There is a finite subset of return lengths \(b_1,\ldots,b_s\)
with gcd one: the successive gcds along an enumeration of the
return lengths are positive integers that eventually stabilize.
The nonnegative combinations of the \(b_\ell\) cover all residue
classes modulo \(b_1\). Indeed, their residue set is closed under
addition, contains zero, and is a subgroup because in a finite
additive group each element has an inverse equal to a nonnegative
multiple of itself. By the gcd condition the subgroup is the
whole group. Choose one nonnegative combination in each residue
class. Every integer larger than all these representatives is
a chosen combination plus a nonnegative multiple of \(b_1\).
Thus all sufficiently large lengths are return lengths at \(a\).
For any \(i,j\), adjoining fixed paths \(i\to a\) and \(a\to j\)
shows \((P^n)_{ij}>0\) for all sufficiently large \(n\).

The independent product chain on \(S\times S\) has transition
matrix
\[
 K_{(i,j),(u,v)}=P_{iu}P_{jv}.
\]
Its \(n\)-step transitions are the products of the two
\(n\)-step probabilities, by the path law or by induction.
The preceding positivity therefore proves irreducibility.
It has stationary probability \(\pi\otimes\pi\), since the two
stationarity sums factor. By Section~\ref{p207:W14} it is positive
recurrent, and Section~\ref{p207:W12} shows it hits the state
\((a,a)\) almost surely from any start. In particular it hits
the diagonal almost surely.

Construct a new pair chain \((U_n,V_n)\) with initial law
\(\delta_i\otimes\pi\). Off the diagonal use \(K\).
At \((j,j)\), draw one common next state from \(P(j,\cdot)\).
Each coordinate retains its original conditional transition law:
both the product kernel and the common draw have the required
one-coordinate marginals. Hence \(U_n\) has law \(P^n(i,\cdot)\)
and \(V_n\) has law \(\pi\).
Let \(T\) be the first time of equality. Its distribution is
the same as in the independent product chain with the same
start, because all path probabilities up to the first equality
use only transitions out of off-diagonal states, apart from a
possible equality at time zero. Thus \(T<\infty\) almost surely.
After \(T\) the coordinates remain equal. The coupling bound
from Section~\ref{p207:W15} gives distance at most
\(\mathbb P(T>n)\to0\).
Finally,
\(d_{\mathrm{TV}}(\lambda P^n,\pi)
\le\sum_i\lambda_i d_{\mathrm{TV}}(P^n(i,\cdot),\pi)\).
The distances are bounded by one, so dominated convergence
for this sum proves the general statement.
\end{proof}

Compare \cite[Theorems 1.8.3--1.8.5]{Norris} and
\cite[Section 5.6]{Durrett}.

\section{Cyclic classes and Ces\`aro convergence}\label{p207:W17}

\begin{theorem}
Let \(P\) be irreducible and positive recurrent with period \(d\).
There is a partition \(S=D_0\sqcup\cdots\sqcup D_{d-1}\) into nonempty
classes such that a one-step transition from \(D_\ell\) is
supported on \(D_{\ell+1}\), with subscripts taken modulo \(d\).
Each class has stationary mass \(1/d\). If \(i\in D_\ell\) and
\(0\le r<d\), then
\[
 P^{md+r}(i,\cdot)\ \longrightarrow\
       d\,\pi(\,\cdot\cap D_{\ell+r})
 \quad\text{in total variation as }m\to\infty.
\]
For every initial probability \(\lambda\),
\[
 \frac1n\sum_{k=0}^{n-1}\lambda P^k
 \ \longrightarrow\ \pi
 \quad\text{in total variation}.
\]
\end{theorem}
\begin{proof}
Fix \(a\in S\). The period \(d\) is positive since \(a\) is recurrent.
For a state \(i\), take any positive-probability path from \(a\) to
\(i\), and assign \(i\) the residue of its length modulo \(d\).
Two such paths to \(i\) can be followed by a common path from \(i\)
back to \(a\); the resulting return lengths differ by the original
difference and are divisible by \(d\). Length zero also counts as
a multiple of \(d\), so the rule includes \(a\).
It is thus well-defined, and appending one edge advances the residue
by one. Every residue occurs along the initial \(d\) positions of
a positive-probability closed path at \(a\), whose length is at least
\(d\).

Write \(\pi_\ell=d\,\pi(\cdot\cap D_\ell)\). Since \(P\) advances
the residue, stationarity gives
\(\pi(D_{\ell+1})=\pi(D_\ell)\), so all these masses are \(1/d\).
In fact the full measure on \(D_{\ell+1}\) comes from \(D_\ell\),
so \(\pi_\ell P=\pi_{\ell+1}\) and \(\pi_\ell P^d=\pi_\ell\).
The restriction of \(P^d\) to \(D_\ell\) is irreducible:
any positive-probability path connecting two states in the same
class has length divisible by \(d\).
Its positive return lengths are exactly the original lengths
divided by \(d\), so its period is one. By the invariant
probability characterization in Section~\ref{p207:W14}, it is
positive recurrent, and by Section~\ref{p207:W16} its transitions
tend to \(\pi_\ell\) in total variation. Applying \(P^r\) and
using contraction proves the residue limit.

For a point start, the distance between the Ces\`aro average
and the average of the corresponding \(\pi_{\ell+k}\) is at
most the average of the termwise distances. Those distances
tend to zero by the finitely many residue limits, so their
averages tend to zero as well. Each residue occurs with limiting
frequency \(1/d\), and \(\sum_\ell\pi_\ell/d=\pi\).
For a general start \(\lambda\), average these bounds over point
starts, using the bound one and dominated convergence.
\end{proof}

If \(d>1\), each point-start time-\(n\) law is carried by a single
cyclic class, so its distance from \(\pi\) is at least \(1-1/d\)
by testing that class. The two-state chain alternating between its
states has precisely this behavior. With a stationary initial
distribution even a periodic chain keeps all marginals stationary.
See \cite[Sections 1.7--1.8]{Norris} and
\cite[Equation (5.6.1)]{Durrett}.

\section{Geometric mixing for finite chains}\label{p207:W18}

\begin{proposition}[Finite minorization bound]
Let \(P\) be finite, irreducible and aperiodic. There is an \(m\ge1\)
such that every entry of \(P^m\) is positive. Put
\(c_j=\min_i(P^m)_{ij}\) and \(\epsilon=\sum_jc_j\in(0,1]\).
For any two initial laws,
\[
 d_{\mathrm{TV}}(\mu P^n,\nu P^n)
 \le (1-\epsilon)^{\lfloor n/m\rfloor}
       d_{\mathrm{TV}}(\mu,\nu),
\]
where a power with exponent zero is one.
\end{proposition}
\begin{proof}
For every pair, positivity holds at all sufficiently large times
by the argument in Section~\ref{p207:W16}. Taking a common upper
bound over finitely many pairs gives \(m\).
The sum \(\epsilon\) is positive and at most any row sum, namely one.
If \(0<\epsilon<1\), let
\[
 \rho_j=c_j/\epsilon,\qquad
 Q_{ij}=\frac{(P^m)_{ij}-c_j}{1-\epsilon}.
\]
Then \(Q\) is stochastic, and for any probability \(\mu\),
\(\mu P^m=\epsilon\rho+(1-\epsilon)\mu Q\).
The common part cancels in the difference of two laws, and
contraction under \(Q\) gives a factor \(1-\epsilon\).
Iterate over blocks and use contraction for the remaining steps.
If \(\epsilon=1\), every row of \(P^m\) equals \(c\), since a
nonnegative coordinatewise difference of sum zero vanishes.
Then the two laws coincide after one block, which is the same bound.
\end{proof}

With \(\nu=\pi\) and \(0<\delta<1\), it suffices to run
\(m\lceil \log(1/\delta)/(-\log(1-\epsilon))\rceil\) steps for
distance at most \(\delta\) when \(\epsilon<1\).
For \(\epsilon=1\), a single block gives exact stationarity.
This is the common-mass method behind the coupling estimates of
\cite[Section 5.6]{Durrett}.

\section{Excursions and the trajectory ergodic theorem}
\label{p207:W19}

\begin{theorem}
Let \(P\) be irreducible and positive recurrent, with stationary
probability \(\pi\). For every real \(f\) on \(S\) with
\(\sum_j\pi_j|f(j)|<\infty\), and every initial law \(\lambda\),
\[
 \frac1n\sum_{k=0}^{n-1}f(X_k)
 \ \longrightarrow\ \sum_j\pi_j f(j)
 \qquad\text{almost surely}.
\]
\end{theorem}
\begin{proof}
Fix a reference state \(i\).
Let \(R_0=T_{\{i\}}\) and let
\(R_{s+1}=\inf\{n>R_s:X_n=i\}\).
By Sections~\ref{p207:W10} and \ref{p207:W12}, all these
stopping times are finite almost surely from every point start,
and hence from any \(\lambda\).
For \(s\ge0\), on the event \(R_s,R_{s+1}<\infty\) set
\[
 L_s=R_{s+1}-R_s,\qquad
 B_s=\sum_{R_s\le k<R_{s+1}}f(X_k),\qquad
 C_s=\sum_{R_s\le k<R_{s+1}}|f(X_k)|.
\]
Each vector \((L_s,B_s,C_s)\) is a measurable function of the
shifted path after \(R_s\), determined by its first positive
return to \(i\). Assign the entire vector the value \((0,0,0)\)
on the null event that \(R_s=\infty\) or \(R_{s+1}=\infty\).
The previous vectors are measurable
in \(\mathcal F_{R_s}\), because on \(\{R_s=n\}\) they are
functions of the prefix through time \(n\). Thus the strong Markov
property gives these vectors a common conditional law given
\(\mathcal F_{R_s}\), the law of one excursion under
\(\mathbb P_i\), independent of that past. Iterating proves that
the sequence of vectors is independent and identically distributed.

With \(\gamma_j=\pi_j/\pi_i\) from Section~\ref{p207:W14}, Tonelli
gives \(\mathbb E C_s=\sum_j\gamma_j|f(j)|<\infty\).
Absolute integrability allows signed summation for \(B_s\),
and the return-time formula gives
\[
 \mu=\mathbb E L_s=1/\pi_i,\qquad
 \mathbb E B_s=\sum_j\gamma_j f(j).
\]
Apply the integrable independent strong law
\cite[Theorem \texttt{thm:etemadi}]{P105} to \(L_s\), \(B_s\),
and \(C_s\). On a common event of probability one,
\(R_m/m=(R_0+\sum_{s<m}L_s)/m\to\mu>0\) and
\(m^{-1}\sum_{s<m}B_s\to\mathbb E B_0\).
Also \(C_m/m\to0\), because
\[
 \frac{C_m}{m}
  =\frac{m+1}{m}\left(\frac1{m+1}\sum_{s\le m}C_s\right)
     -\frac1m\sum_{s<m}C_s
\]
and the two normalized sums have the same limit \(\mathbb E C_0\).

For \(n\ge R_0\), choose \(m=m(n)\) with \(R_m\le n<R_{m+1}\).
The times are strictly increasing, so \(m(n)\to\infty\), and
the inequalities together with \(R_m/m\to\mu\) give
\(n/m(n)\to\mu\). The initial segment before \(R_0\) contributes
a finite sum, which vanishes after division by \(n\).
The incomplete last excursion contributes in absolute value at
most \(C_m\), and \(C_m/n\to0\).
The contribution from the \(m\) full excursions is
\((m/n)(m^{-1}\sum_{s<m} B_s)\), tending to
\(\mathbb E B_0/\mu=\sum_j\pi_j f(j)\).
\end{proof}

The regenerative proof follows the time-average approach in
\cite[Section 1.10]{Norris} and \cite[Section 5.7]{Durrett}.
The absolute occupation bound and unfinished-excursion estimate
give the stated \(L^1(\pi)\) scope. No aperiodicity hypothesis
enters the proof.
\chapter[Random Walks]{Random Walks on Intervals, Lattices, and Graphs}

We use \(\mathbb P_i\) and \(\mathbb E_i\) for the law and expectation
with starting point \(i\). In this chapter
\(T_A=\inf\{n\geq 0:X_n\in A\}\), with \(\inf\varnothing=\infty\).
For a single point \(a\) write \(T_a=T_{\{a\}}\).
The first positive return to the starting point is
\(T_i^+=\inf\{n\geq 1:X_n=i\}\).

\section{Gambler's ruin}\label{p207:W20}

\begin{theorem}\label{p207:thm:ruin}
Let \(N\geq 1\) and \(0<p<1\), and put \(q=1-p\).
On \(\{0,\ldots,N\}\), move right with probability \(p\) and left with
probability \(q\) at each interior point, and absorb at both endpoints.
Then \(T=T_{\{0,N\}}\) is almost surely finite and has finite mean.
For every \(i\in\{0,\ldots,N\}\), the upper exit probability is
\[
 h_i=\mathbb P_i(T_N<T_0)=
 \begin{cases}
  i/N, & p=q,\\[2pt]
  \displaystyle\frac{1-(q/p)^i}{1-(q/p)^N}, & p\ne q .
 \end{cases}
\]
\end{theorem}

\begin{proof}
Realize the moves by independent right or left choices even after the
walk has reached the boundary, where further choices leave it unchanged.
Conditional on survival to time \(kN\), the next \(N\) choices are all
left with probability \(q^N\). Such choices force absorption by
\((k+1)N\), from every interior starting point. Consequently
\[
 \mathbb P_i(T>kN)\leq (1-q^N)^k\qquad(k\geq 0).
\]
The tail probabilities tend to zero. Their sum over all integer times
is at most \(N\sum_{k\geq 0}(1-q^N)^k\). By summing the indicators
\(1_{\{T>m\}}\), this tail sum equals \(\mathbb E_i T\).

The endpoint values are \(h_0=0\) and \(h_N=1\).
At an interior point, conditioning on the first step gives
\(h_i=p h_{i+1}+q h_{i-1}\). For \(d_i=h_i-h_{i-1}\) this becomes
\(d_{i+1}=(q/p)d_i\) for \(1\leq i<N\). Thus
\[
 h_i=d_1\sum_{j=0}^{i-1}(q/p)^j,\qquad
 1=d_1\sum_{j=0}^{N-1}(q/p)^j.
\]
The second sum is positive, so the solution is unique. Evaluating the
two sums when \(p=q\) and when \(p\ne q\) gives the formulas.
For \(N=1\) there are only the prescribed endpoint values, and the
same formulas hold.
\end{proof}

For textbook treatments of the first-step equations and ruin,
see \cite[\S\S1.3 and 1.6]{Norris} and \cite[\S4.8.1]{Durrett}.

\section{Absorption times and biased walks}\label{p207:W21}

\begin{theorem}\label{p207:thm:ruin-time}
Under the assumptions of Theorem~\ref{p207:thm:ruin}, the mean absorption
time is
\[
 \mathbb E_i T=
 \begin{cases}
  i(N-i), & p=q,\\[2pt]
  (N h_i-i)/(p-q), & p\ne q.
 \end{cases}
\]
For the same independent right or left increments on the whole integer
line, the probability of reaching zero from \(i>0\) is
\[
 \mathbb P_i(T_0<\infty)=
 \begin{cases}
  (q/p)^i, & p>q,\\
  1, & p\leq q .
 \end{cases}
\]
For \(i<0\) the corresponding probability is \((p/q)^{-i}\) if \(p<q\),
and \(1\) if \(p\geq q\). Starting from zero, the probability of a
positive return to zero is \(2\min(p,q)\). In the symmetric case,
\(\mathbb E_1 T_0=\mathbb E_{-1} T_0=\mathbb E_0 T_0^+=\infty\).
\end{theorem}

\begin{proof}
Write \(g_i=\mathbb E_i T\), which is finite by
Theorem~\ref{p207:thm:ruin}. The first-step equation is
\(g_i=1+p g_{i+1}+q g_{i-1}\) at interior points, and
\(g_0=g_N=0\). With \(L f(i)=p f(i+1)+q f(i-1)-f(i)\) on the interior,
it suffices to solve \(L g=-1\) with these boundary values.
The difference of two solutions satisfies \(L f=0\) and vanishes at
the endpoints, so the difference recurrence in the preceding proof
gives \(f=0\). If \(p=q\), the centered second difference of
\(i(N-i)\) is \(-2\), giving \(L g=-1\).
If \(p\ne q\), \(Lh=0\) and \(L(i\mapsto i)=p-q\). The displayed
expression then solves the equation and takes value zero at each
endpoint.

Fix \(i>0\). For \(N>i\) the events \(\{T_0<T_N\}\) for the walk on the
line increase as \(N\) increases, since a nearest-neighbor path that
hits \(N+1\) must have hit \(N\). Their union is \(\{T_0<\infty\}\),
because every finite path before the first hit of zero has a finite
maximum. Their probabilities equal \(1-h_i\) by
Theorem~\ref{p207:thm:ruin}. In the symmetric case \(1-h_i=1-i/N\)
tends to \(1\). With \(r=q/p\ne 1\), the value is
\[
 1-h_i=\frac{r^i-r^N}{1-r^N}.
\]
It tends to \(r^i\) if \(r<1\) and to \(1\) if \(r>1\).
Reflection of the walk exchanges \(p\) and \(q\) and yields the
formula from negative points. From zero, condition on the first move.
When \(p>q\), the positive return probability is \(p(q/p)+q=2q\).
When \(p<q\) it is \(p+q(p/q)=2p\), and in the symmetric case it is \(1\).

When \(p=q\), \(T_{\{0,N\}}\leq T_0\) from \(1\) and the mean of the left
side equals \(N-1\) for every \(N>1\). Thus
\(\mathbb E_1T_0=\infty\). Reflection gives the same result from \(-1\),
and conditioning on the first step gives
\(\mathbb E_0 T_0^+=\infty\).
\end{proof}

\section{Pólya's lattice recurrence criterion}\label{p207:W22}

Fix \(d\geq 1\). The simple symmetric walk on \(\mathbb Z^d\) is
\(X_n=X_0+\xi_1+\cdots+\xi_n\), where the \(\xi_i\) are independent and
each of the \(2d\) points \(\pm e_j\) has probability \(1/(2d)\).
Write \(\mathbb T^d=(\mathbb R/(2\pi\mathbb Z))^d\) and denote by \(m\)
its Lebesgue probability measure, that is, ordinary Lebesgue measure
on a fundamental cube divided by \((2\pi)^d\).
The Fourier argument below follows the lattice-walk framework in
\cite[\S5.4]{Durrett} and \cite[\S3.8]{Knill}.

\begin{lemma}\label{p207:lem:lattice-fourier}
Set
\[
 \phi(t)=\frac1d\sum_{j=1}^d\cos t_j .
\]
The origin return probability at time \(n\) is
\(\displaystyle P^n_{00}=\int_{\mathbb T^d}\phi(t)^n\,dm(t)\).
It is zero for odd \(n\). The points with \(\lvert\phi\rvert=1\) are
precisely \(0\) and \(\boldsymbol\pi=(\pi,\ldots,\pi)\) in \(\mathbb T^d\).
For each of these two points \(t_*\) there is a local Euclidean
coordinate \(u\) centered at \(t_*\) and positive constants \(a,b,\eta\)
such that, whenever \(\lVert u\rVert\leq\eta\),
\[
 a\lVert u\rVert^2\leq 1-\lvert\phi(t_*+u)\rvert
                  \leq b\lVert u\rVert^2 .
\]
\end{lemma}

\begin{proof}
For \(z\in\mathbb Z^d\), the integral of
\(e^{\mathrm i\langle t,z\rangle}\) equals \(1\) for \(z=0\) and \(0\)
otherwise. This follows by taking the product of the one-dimensional
integrals of \(e^{\mathrm i t_j z_j}\) over a period. In particular,
when the walk starts at the origin, integrating the characteristic
function of \(X_n\) retains precisely its mass at zero. This
interchange uses only a finite sum at each time. The independence of
the increments gives the \(n\)th power of their one-step characteristic
function, which is \(\phi\). The sum of the integer coordinates
changes parity at every step, so the walk cannot return to zero at an
odd time.

The average of the \(d\) cosines is \(1\) only if all of them are \(1\).
It is \(-1\) only if all of them are \(-1\). These two cases give the
specified points on the torus. In a ball about \(0\) of sufficiently
small radius, \(\phi(u)>0\). The inequalities
\(a_0 s^2\leq 1-\cos s\leq b_0 s^2\) for small \(\lvert s\rvert\)
follow from \(\lim_{s\to0}(1-\cos s)/s^2=1/2\), for fixed positive
\(a_0,b_0\). Summing over the coordinates gives the claimed bounds
at \(0\), with constants \(a_0/d,b_0/d\).
The identity \(\phi(\boldsymbol\pi+u)=-\phi(u)\) gives the same
bounds in the translated chart.
\end{proof}

\begin{theorem}[Pólya's classification]\label{p207:thm:polya}
The simple symmetric walk on \(\mathbb Z^d\) is recurrent for \(d=1,2\)
and transient for \(d\geq3\).
\end{theorem}

\begin{proof}
Translations preserve the increment law, so all points have the
same return probabilities. By the Green-series criterion of
Section~\ref{p207:W10}, it suffices to study their sum at the origin.
The formula and parity statement in
Lemma~\ref{p207:lem:lattice-fourier} and monotone convergence give
\[
 \sum_{n\geq 0}P^n_{00}
 =\sum_{n\geq 0}\int_{\mathbb T^d}\phi(t)^{2n}\,dm(t)
 =\int_{\mathbb T^d}\frac{1}{1-\phi(t)^2}\,dm(t),
\]
where the integrand takes the extended value \(+\infty\) at the two
specified points. Choose disjoint chart balls of radius \(\eta>0\)
around those points, shrinking the radius if necessary. Their
complement is compact and \(\lvert\phi\rvert<1\) there, so the
integrand is bounded on the complement. In each ball,
\(1-\phi^2=(1-\lvert\phi\rvert)(1+\lvert\phi\rvert)\), and the second
factor lies between \(1\) and \(2\). Hence the local integrand is
bounded above and below by positive constants times
\(\lVert u\rVert^{-2}\).

To decide the local integral, use the annuli
\(\eta2^{-(j+1)}<\lVert u\rVert\leq \eta2^{-j}\), \(j\geq0\).
Each has ordinary volume
\(\omega_d\eta^d(1-2^{-d})2^{-jd}\), where \(\omega_d\) is the volume of
the unit ball. On the annulus, \(\lVert u\rVert^{-2}\) is bounded
between \(\eta^{-2}2^{2j}\) and \(4\eta^{-2}2^{2j}\). Summing the
resulting upper and lower bounds shows that the local integral is
finite exactly when \(\sum_{j\geq0}2^{-j(d-2)}\) converges, that is,
when \(d>2\). The Green criterion completes the classification.
\end{proof}

\section{Even return probabilities}\label{p207:W23}

\begin{theorem}\label{p207:thm:lattice-return-rate}
For each \(d\geq1\) there are finite constants \(0<c_d\leq C_d\) such
that
\[
 c_d\,n^{-d/2}\leq P^{2n}_{00}\leq C_d\,n^{-d/2}
 \qquad(n\geq1)
\]
for the walk of Section~\ref{p207:W22}.
\end{theorem}

\begin{proof}
Use the disjoint chart balls of radius \(\eta\) from the previous
proof. In either ball, Lemma~\ref{p207:lem:lattice-fourier} gives
\[
 \lvert\phi(t_*+u)\rvert^{2n}
 \leq (1-a\lVert u\rVert^2)^{2n}
 \leq \exp(-2an\lVert u\rVert^2).
\]
The first base on the right is nonnegative by the same local estimate.
The integral of the last expression over either ball is bounded by
the integral over \(\mathbb R^d\), times \((2\pi)^{-d}\).
The change of variables \(v=\sqrt n\,u\) makes this bound a constant
times \(n^{-d/2}\). The integral of \(e^{-2a\lVert v\rVert^2}\) is
finite by its product form and the one-dimensional bound
\(e^{-2a s^2}\leq e^{-2a\lvert s\rvert}\) when \(\lvert s\rvert\geq1\).
On the complement of the two balls there is a uniform bound
\(\lvert\phi\rvert\leq\theta<1\). If \(\theta=0\) this region contributes
zero for \(n\geq1\). If \(0<\theta<1\), \(n^{d/2}\theta^{2n}\) is
bounded, so the complement satisfies the same order bound. Integrating
the nonnegative even power proves the upper estimate.

Choose \(s>0\) with \(s\leq\eta\) and \(b s^2\leq1/4\). On the ball
\(\lVert u\rVert\leq s/\sqrt n\) about zero, the other local inequality
gives \(\lvert\phi(u)\rvert\geq 1-1/(4n)\). Bernoulli's inequality
\((1-x)^k\geq1-kx\) for \(0\leq x\leq1\), which follows by induction
on \(k\), gives
\[
 \lvert\phi(u)\rvert^{2n}\geq \bigl(1-1/(4n)\bigr)^{2n}\geq \tfrac12 .
\]
The ball has \(m\)-measure
\((2\pi)^{-d}\omega_d s^d n^{-d/2}\), proving the lower estimate.
\end{proof}

\section{Walks on finite connected graphs}\label{p207:W24}

\begin{theorem}\label{p207:thm:graph-walk}
Let \(G=(V,E)\) be a finite connected undirected simple graph with
at least two vertices. The kernel that chooses a neighbor uniformly is
\[
 P_{uv}=
 \begin{cases}
  1/\deg(u), & \{u,v\}\in E,\\
  0, & \{u,v\}\notin E.
 \end{cases}
\]
The kernel is reversible with unique stationary probability
\[
 \pi(u)=\frac{\deg(u)}{2\#E}.
\]
The period is \(2\) when \(G\) is bipartite and \(1\) when \(G\) is
nonbipartite. In the second case the walk converges to \(\pi\) in
total variation from every starting law.
For any \(0<\alpha<1\), the kernel \(\alpha I+(1-\alpha)P\)
has the same stationary probability and is aperiodic, so its laws
converge to \(\pi\) in total variation from every starting law.
\end{theorem}

\begin{proof}
Connectivity and the size assumption ensure that all degrees are
positive and that \(P\) is irreducible. Since each edge contributes
twice to the sum of degrees, the displayed formula defines a
probability. Across an edge \(\{u,v\}\), both \(\pi(u)P_{uv}\)
and \(\pi(v)P_{vu}\) equal \(1/(2\#E)\). For all other pairs both
are zero. Thus detailed balance holds. By the results of
Sections~\ref{p207:W13} and~\ref{p207:W14}, \(\pi\) is reversible
and is the unique stationary probability.

Every vertex has a two-step return through one of its neighbors.
In a bipartite graph every closed walk has even length, since each
step exchanges the two parts. Thus the period is \(2\).
Conversely, if all closed walks have even length, fix a root and
give each vertex the parity of the length of a path from the root.
Two paths to the same vertex have the same parity, since one
followed by the reverse of the other is a closed walk. An edge
changes parity, yielding a bipartition. Therefore a nonbipartite
graph has an odd closed walk. From any vertex one can walk to a
vertex of that closed walk, traverse it, then reverse the access
path. The resulting return has odd length and positive
probability. Together with a length-two return this gives period
\(1\). The total variation convergence claim now follows from
Section~\ref{p207:W16}. In the bipartite case the ordinary laws from
a single vertex alternate parts, and
Section~\ref{p207:W17} gives the cyclic and Ces\`aro limits.

For \(\alpha I+(1-\alpha)P\), the same detailed balance identity
holds by linearity. The kernel remains irreducible because every
original edge keeps positive transition probability. Its
one-step return at every vertex has positive probability, so
Section~\ref{p207:W16} applies again.
\end{proof}

The detailed-balance and period arguments use the finite-chain
framework of \cite[\S\S1.8--1.9]{Norris} and
\cite[\S\S5.5--5.6]{Durrett}.

\section{Dirichlet forms and spectral mixing}
\label{p207:W25}

Let \(P\) be an irreducible reversible kernel on a finite state space
\(S\), and let \(\pi\) denote its stationary probability, which is
strictly positive by Section~\ref{p207:W13}. Write
\(\langle f,g\rangle_\pi=\sum_{x\in S}\pi(x)f(x)g(x)\) for the real
\(\ell^2(\pi)\) inner product, and
\(\lVert f\rVert_{2,\pi}^2=\langle f,f\rangle_\pi\).
The action on functions is \(P f(x)=\sum_y P_{xy}f(y)\).

\begin{theorem}\label{p207:thm:reversible-form}
The operator \(P\) is self-adjoint and has norm at most \(1\) on
\(\ell^2(\pi)\). For all real \(f,g\), its Dirichlet form is
\[
 \mathcal E(f,g)=\langle f,(I-P)g\rangle_\pi
   =\frac12\sum_{x,y\in S}
       \pi(x)P_{xy}\bigl(f(x)-f(y)\bigr)\bigl(g(x)-g(y)\bigr).
\]
The nullspace of \(I-P\) consists of the constant functions.
\end{theorem}

\begin{proof}
By detailed balance, switching the indices \(x,y\) in the finite sum
for \(\langle f,Pg\rangle_\pi\) turns it into
\(\langle Pf,g\rangle_\pi\). For each \(x\), the weighted Cauchy--Schwarz
inequality gives \((Pf(x))^2\leq\sum_yP_{xy}f(y)^2\). Multiplying by
\(\pi(x)\), summing over \(x\), and using stationarity gives
\(\lVert Pf\rVert_{2,\pi}^2\leq \lVert f\rVert_{2,\pi}^2\).
To expand the Dirichlet sum, the terms involving \(f(x)g(x)\)
use the row sums \(1\), and those involving \(f(y)g(y)\)
use the column stationarity sums \(\sum_x\pi(x)P_{xy}=\pi(y)\).
The two mixed terms agree by detailed balance. This gives the
formula. In particular, \(\mathcal E(f,f)=0\) only if
\(f(x)=f(y)\) whenever \(P_{xy}>0\). Irreducibility then makes \(f\)
constant on \(S\). Constant functions have \(P f=f\), proving
the nullspace assertion.
\end{proof}

\begin{theorem}\label{p207:thm:spectral-tv}
Suppose \(S\) has at least two states, and let \(r\) be the maximum
absolute value of the eigenvalues of \(P\) on the mean-zero subspace
\[
 H_0=\Bigl\{f:\sum_x\pi(x)f(x)=0\Bigr\}.
\]
Then for every \(i\in S\) and every \(n\geq0\),
\[
 d_{\mathrm{TV}}\bigl(P^n(i,\cdot),\pi\bigr)
 \leq \tfrac12\sqrt{\pi(i)^{-1}-1}\,r^n ,
\]
with \(r^0=1\). In the half-lazy kernel \((I+P)/2\), every eigenvalue
is nonnegative and the spectral radius on \(H_0\) is strictly less
than \(1\).
\end{theorem}

\begin{proof}
The space \(H_0\) is invariant under \(P\) by stationarity. By the
finite-dimensional real spectral theorem
\cite[The Real Spectral Theorem]{D05}, the self-adjoint restriction
has an orthonormal eigenbasis. The eigenvalue bound in the preceding
theorem places all eigenvalues in \([-1,1]\). Expanding in the
eigenbasis gives
\(\lVert P^n f\rVert_{2,\pi}\leq r^n\lVert f\rVert_{2,\pi}\)
for all \(f\in H_0\).

For \(\mu_n=P^n(i,\cdot)\), define the density
\(\rho_n(x)=\mu_n(x)/\pi(x)\). Detailed balance gives
\[
 \rho_{n+1}(x)=\sum_y \frac{\pi(y)}{\pi(x)} P_{yx}\rho_n(y)
             =\sum_y P_{xy}\rho_n(y)=P\rho_n(x).
\]
Consequently \(\rho_n-1=P^n(\rho_0-1)\), with \(\rho_0-1\in H_0\).
Since \(\rho_0(x)=1_{\{i\}}(x)/\pi(i)\), direct summation yields
\(\lVert \rho_0-1\rVert_{2,\pi}^2=\pi(i)^{-1}-1\).
The finite-state total variation formula of Section~\ref{p207:W15}
and Cauchy--Schwarz for the probability \(\pi\) give
\[
 d_{\mathrm{TV}}(\mu_n,\pi)
  =\tfrac12\sum_x\pi(x)\lvert \rho_n(x)-1\rvert
  \leq\tfrac12\lVert \rho_n-1\rVert_{2,\pi}.
\]
This proves the stated bound. The eigenvalue \(\lambda\) becomes
\((1+\lambda)/2\) under half-lazification, so the new eigenvalues
are in \([0,1]\). The preceding nullspace result excludes
\(\lambda=1\) on \(H_0\). Finiteness of the spectrum makes the
largest transformed eigenvalue on \(H_0\) strictly less than \(1\).
\end{proof}

\begin{theorem}[Dirichlet principle]\label{p207:thm:dirichlet-principle}
For a nonempty boundary \(B\subseteq S\) and prescribed real boundary
values \(g:B\to\mathbb R\), there is a unique \(h:S\to\mathbb R\)
satisfying \(h=g\) on \(B\) and \(Ph=h\) on \(S\setminus B\).
It is given by \(h(x)=\mathbb E_x[g(X_{T_B})]\) and is the unique
minimizer of \(\mathcal E(f,f)\) among functions \(f\) with \(f=g\)
on \(B\).
\end{theorem}

\begin{proof}
From each point of \(S\setminus B\) choose a finite positive-probability
path to \(B\). If the complement is nonempty, there is a common bound
\(m\geq1\) on path length and a positive common lower bound
\(\varepsilon\) on path probability. By the deterministic-time
Markov property, conditional on the state at a block start and on
survival to that time, the chance of hitting \(B\) in the next
\(m\) steps is at least \(\varepsilon\). Induction gives
\(\mathbb P_x(T_B>km)\leq(1-\varepsilon)^k\), so \(T_B\) is
almost surely finite. If \(B=S\), \(T_B=0\). The bounded function
\(g\) therefore defines the indicated expectation at every point.
At points of \(B\) it equals \(g\), and at other points the first-step
equation gives \(Ph=h\).

If \(f=h+w\) with \(w=0\) on \(B\), self-adjointness gives
\(\mathcal E(h,w)=\mathcal E(w,h)=\langle w,(I-P)h\rangle_\pi=0\).
Each summand vanishes because \(w=0\) on \(B\) and \(h\) is harmonic
on \(S\setminus B\). Hence
\(\mathcal E(f,f)=\mathcal E(h,h)+\mathcal E(w,w)\).
The last term is nonnegative and, if zero, makes \(w\) constant by
the formula in Theorem~\ref{p207:thm:reversible-form}. Since \(B\)
is nonempty, the constant is zero. This proves uniqueness of the
minimizer. Any other function harmonic off \(B\) and equal to \(g\)
on \(B\) gives the same orthogonality argument and so must be that
minimizer.
\end{proof}

For the boundary hitting equations underlying this representation,
see \cite[\S1.3]{Norris}.
\chapter{Maximal Inequalities and Upcrossings}

\section{The finite-horizon weak maximal inequality}\label{p207:W26}

\begin{theorem}[Doob's weak maximal inequality]
Let \(X_0,\ldots,X_n\) be a real submartingale and put
\(Z=\max_{0\le k\le n}X_k\). For every \(\lambda>0\),
\[
 \lambda\mathbb P(Z\ge\lambda)
 \le \mathbb E[X_n1_{\{Z\ge\lambda\}}]
 \le \mathbb E X_n^+.
\]
If \(X\) is nonnegative, the last term is \(\mathbb E X_n\).
\end{theorem}
\begin{proof}
For each \(0\le k\le n\), let \(A_k\) be the event that \(k\)
is the first time with \(X_k\ge\lambda\). It belongs to
\(\mathcal F_k\), since it also tests only the values at times
less than \(k\). The events are disjoint and have union
\(\{Z\ge\lambda\}\). The two-time submartingale property gives
\[
 \lambda\mathbb P(A_k)
 \le \mathbb E[X_k1_{A_k}]
 \le \mathbb E[X_n1_{A_k}].
\]
Sum over \(k\). The final upper bound follows by comparing
\(X_n1_{\{Z\ge\lambda\}}\) pointwise to \(X_n^+\).
\end{proof}

See \cite[Section 14.6]{Williams}, \cite[Theorem 9.16]{Kallenberg},
and \cite[Theorem 4.4.2]{Durrett}.

\section{The finite-horizon strong maximal inequality}\label{p207:W27}

\begin{theorem}[Doob's strong \(L^p\) inequality]
Let \(X_0,\ldots,X_n\) be a nonnegative real submartingale,
\(1<p<\infty\), and \(X_n\in L^p\). For \(Z=\max_{k\le n}X_k\),
\[
     \lVert Z\rVert_p\le \frac{p}{p-1}\lVert X_n\rVert_p .
\]
In particular, a real martingale \(M\) with \(M_n\in L^p\) satisfies
the same bound with \(Z=\max_{k\le n}|M_k|\) and
\(\lVert X_n\rVert_p\) replaced by \(\lVert M_n\rVert_p\).
\end{theorem}
\begin{proof}
For \(k\le n\), \(0\le X_k\le \mathbb E[X_n\mid\mathcal F_k]\).
Conditional Jensen for \(x\mapsto |x|^p\) gives
\(\mathbb E X_k^p\le\mathbb E X_n^p\)
\cite[conditional Jensen theorem]{P105}. Thus \(Z^p\) is integrable,
since \(Z^p\le\sum_{k\le n}X_k^p\). For each sample point,
\(Z^p=\int_0^\infty p\lambda^{p-1}1_{\{Z\ge\lambda\}}\,d\lambda\).
Use Tonelli and the localized inequality in
Section~\ref{p207:W26}, and then Tonelli again, to get
\begin{align*}
 \mathbb E Z^p
 &\le \int_0^\infty p\lambda^{p-2}
                 \mathbb E[X_n1_{\{Z\ge\lambda\}}]\,d\lambda \\
 &= \frac{p}{p-1}\mathbb E[X_n Z^{p-1}]\\
 &\le \frac{p}{p-1}\lVert X_n\rVert_p\lVert Z\rVert_p^{p-1}.
\end{align*}
The integral at zero is finite for each finite \(Z\) because \(p>1\).
If \(\lVert Z\rVert_p=0\) the claim holds, and otherwise division
gives the bound. For a martingale, conditional Jensen for the
absolute value makes \((|M_k|)\) a nonnegative submartingale.
\end{proof}

Compare \cite[Section 14.11]{Williams},
\cite[Theorem 9.17]{Kallenberg}, and \cite[Theorem 4.4.4]{Durrett}.

\section{The finite-horizon upcrossing inequality}\label{p207:W28}

Fix \(a<b\). For a real sequence \(x_0,\ldots,x_n\), an upcrossing
of \([a,b]\) consists of a time at which \(x\le a\) followed by
a strictly later time at which \(x\ge b\). The greedy count
\(U_n(a,b)\) is computed by buying at the first available time
at or below \(a\), selling at the first later time at or above
\(b\), and repeating from the time after the sale. It counts
only completed upcrossings, including a possible sale at \(n\).

\begin{lemma}\label{p207:greedy}
The greedy count equals the largest \(m\) admitting ordered
pairs \(0\le s_1<t_1<s_2<t_2<\cdots<s_m<t_m\le n\) with
\(x_{s_j}\le a\) and \(x_{t_j}\ge b\).
\end{lemma}
\begin{proof}
Greedy times themselves satisfy the conditions. Given any
ordered pairs, the first greedy buy is no later than \(s_1\)
and the first greedy sale is no later than \(t_1\). If the
greedy sale at step \(j\) is no later than \(t_j\), the next
specified buy \(s_{j+1}\) remains available, so the next greedy
buy and sale are no later than \(s_{j+1}\) and \(t_{j+1}\).
Induction proves that the greedy algorithm completes at
least as many pairs as the specified sequence.
\end{proof}

\begin{theorem}[Doob's upcrossing inequality]\label{p207:upcross}
Let \(X\) be an integrable real submartingale. For the count
along \(X_0,\ldots,X_n\),
\begin{align*}
 (b-a)\mathbb E U_n(a,b)
  &\le \mathbb E[(X_n-a)^+]-\mathbb E[(X_0-a)^+]\\
  &\le \mathbb E[(X_n-a)^+]+\mathbb E[(X_0-a)^-].
\end{align*}
\end{theorem}
\begin{proof}
The process \(Y_k=(X_k-a)^+\) is a nonnegative submartingale:
conditional Jensen gives the first comparison to
\((\mathbb E[X_{k+1}\mid\mathcal F_k]-a)^+\), and the increasing
positive-part map and the submartingale inequality give the
second comparison to \(Y_k\). Apply the greedy rule to \(Y\)
with buy threshold \(0\) and sell threshold \(b-a\). These
events are exactly the buy and sell events for \(X\).
Let \(H_k=1\) if the rule holds a unit after processing time \(k\)
and \(H_k=0\) otherwise. Its state depends measurably on the
values through time \(k\), so \(H\) is predictable in the sense
of Section~\ref{p207:W03}. The count is also \(\mathcal F_n\)-measurable,
being determined by finitely many measurable tests.

The gain \(\sum_{k<n}H_k(Y_{k+1}-Y_k)\) telescopes on each
holding interval. Every completed interval earns at least
\(b-a\) because its buy value is zero, and an interval open
at the cutoff earns \(Y_n\ge0\). A buy made at \(n\) contributes
zero. Thus the gain is at least \((b-a)U_n(a,b)\).
The gain of \(1-H\) has nonnegative expectation by the
predictable transform inequality, and the two gains sum to
\(Y_n-Y_0\). This gives the first bound. The second follows
from \(-\mathbb E[(X_0-a)^+]\le \mathbb E[(X_0-a)^-]\).
\end{proof}

See \cite[Sections 11.2--11.5]{Williams} and
\cite[Lemma 9.18]{Kallenberg}. The endpoint convention also
fixes how to compare this statement to a strict cutoff in
the fixed Mathlib implementation \cite{Mathlib}.

\section{Azuma--Hoeffding for bounded increments}\label{p207:W29}

\begin{theorem}
Let \(M_0,\ldots,M_n\) be a real martingale and let
\(a_k\le M_{k+1}-M_k\le b_k\) almost surely for deterministic
finite bounds \(a_k,b_k\), \(0\le k<n\).
Set \(V=\sum_{k<n}(b_k-a_k)^2\). If \(V>0\) and \(t>0\), then
\[
 \max\{\mathbb P(M_n-M_0\ge t),\,\mathbb P(M_n-M_0\le-t)\}
       \le e^{-2t^2/V}.
\]
Consequently the two-sided probability is bounded by
\(2e^{-2t^2/V}\). If \(V=0\), both one-sided probabilities are zero
for every \(t>0\).
\end{theorem}
\begin{proof}
Fix an increment \(D=M_{k+1}-M_k\) with bounds \(a,b\).
Its conditional mean given \(\mathcal F_k\) is zero,
so \(a\le0\le b\), and \(a=b\) forces \(D=0\) almost surely.
For \(a<b\), convexity of the exponential gives the chord bound
\[
 e^{sD}\le \frac{b-D}{b-a}e^{sa}+\frac{D-a}{b-a}e^{sb}.
\]
Taking conditional expectations removes \(D\) from the weights.
The bound is the moment-generating function of a two-point
variable with mass \(b/(b-a)\) at \(a\) and mass \(-a/(b-a)\)
at \(b\), and mean zero. To see its logarithm \(g(s)\) is at most
\(s^2(b-a)^2/8\), differentiate the two-term sum: if \(u_s\) is
the mass at \(b\) after weighting by the exponentials and
renormalizing, then \(g''(s)=u_s(1-u_s)(b-a)^2\le(b-a)^2/4\).
The same formula gives zero variance when one mass vanishes.
Now \(g(0)=g'(0)=0\), so Taylor's integral formula bounds
\(g(s)\) as stated for either sign of \(s\).
This is the bounded-variable logarithmic moment estimate of
\cite[W25]{P205} specialized to the chord law.

Each partial sum \(M_k-M_0\) is \(\mathcal F_k\)-measurable
and bounded, since its increments are bounded.
The tower property and pull-out rule applied successively give
\(\mathbb E e^{s(M_n-M_0)}\le e^{s^2 V/8}\).
For \(s>0\), Markov's inequality bounds the upper tail by
\(e^{-st+s^2V/8}\), whose value at \(s=4t/V\) is \(e^{-2t^2/V}\).
Apply the same argument to \(-M\) for the other tail and use the
union bound. When \(V=0\) every increment is zero.
\end{proof}

For the classical martingale bound, see \cite{Azuma}.
\chapter{Uniform Integrability and Martingale Convergence}
\chaptermark{Uniform Integrability and Convergence}

Throughout this chapter \((\Omega,\mathcal F,\mathbb P)\) is a
probability space. Integrability of a Banach-valued random variable is
Bochner integrability. Such variables have strongly measurable
representatives, meaning almost everywhere limits of measurable simple
functions. One can take the limit on the measurable set of convergence
and set the value to zero elsewhere. All the martingales and
submartingales in the convergence theorems below are real-valued.

\section{Uniform integrability and the doubling example}
\label{p207:W30}

\begin{definition}[Uniform integrability]
A family \((X_i)_{i\in I}\) of integrable random variables with values in
a real or complex Banach space is uniformly integrable if
\[
 \lim_{R\to\infty}\sup_{i\in I}
 \mathbb E\big[\lVert X_i\rVert
                \mathbf 1_{\{\lVert X_i\rVert>R\}}\big]=0.
\]
For real values the norm is the absolute value. We take suprema of
empty families of nonnegative bounds to be zero.
\end{definition}

\begin{proposition}
The definition is unchanged if \(>R\) is replaced by \(\ge R\).
Every uniformly integrable family is bounded in \(L^1\).
\end{proposition}

\begin{proof}
For \(R>0\), the indicator of \(\{\lVert X_i\rVert\ge R\}\) is bounded by
that of \(\{\lVert X_i\rVert>R/2\}\) and bounds that of
\(\{\lVert X_i\rVert>R\}\). For the second claim choose \(R\) large
enough that the supremum of the tails is at most one. Splitting at
\(R\) gives \(\mathbb E\lVert X_i\rVert\le R+1\) for all \(i\).
\end{proof}

\begin{example}[The doubling martingale]
Let \(\mathcal F_n\) reveal the first \(n\) independent fair coin tosses
and let \(A_n\) be the event that all these tosses are heads.
Take \(A_0=\Omega\) and \(M_n=2^n\mathbf 1_{A_n}\).
Since
\(\mathbb E[\mathbf 1_{A_{n+1}}\mid\mathcal F_n]
=\mathbf 1_{A_n}/2\), the process is a martingale.
It has \(\mathbb E|M_n|=1\) for every \(n\).
The probability of all tosses being heads is
\(\lim_n2^{-n}=0\), so \(M_n\to0\) almost surely.
For each \(R\), choosing \(n\) with \(2^n>R\) gives
\(\mathbb E[M_n\mathbf 1_{\{M_n>R\}}]=1\).
The family is therefore neither uniformly integrable nor convergent in
\(L^1\).
\end{example}

\section{Small sets and the de la Vallée--Poussin criterion}
\label{p207:W31}

\begin{theorem}
For an integrable family \((X_i)_{i\in I}\) as in
Section~\ref{p207:W30}, the following conditions are equivalent.
\begin{enumerate}
 \item The family is uniformly integrable.
 \item The family is bounded in \(L^1\), and for every \(\epsilon>0\)
       there is \(\delta>0\) such that
       \(\mathbb E[\lVert X_i\rVert\mathbf 1_A]<\epsilon\)
       for all \(i\) and all measurable \(A\) with
       \(\mathbb P(A)<\delta\).
 \item There is a finite nonnegative, nondecreasing convex function
       \(\Phi\) on \([0,\infty)\) such that
       \[
          \lim_{t\to\infty}\frac{\Phi(t)}{t}=\infty
          \qquad\text{and}\qquad
          \sup_i\mathbb E\Phi(\lVert X_i\rVert)<\infty.
       \]
\end{enumerate}
\end{theorem}

\begin{proof}
Write \(V_i=\lVert X_i\rVert\) and \(C=\sup_i\mathbb EV_i\).
Under uniform integrability the \(L^1\) bound follows from
Section~\ref{p207:W30}, and
\[
 \mathbb E[V_i\mathbf 1_A]
 \le R\,\mathbb P(A)+\mathbb E[V_i\mathbf 1_{\{V_i>R\}}].
\]
First choose \(R>0\) to bound the second term by \(\epsilon/2\),
then choose \(\delta=\epsilon/(2R)\).
Conversely if \(C<\infty\), Markov's inequality
gives \(\mathbb P(V_i>R)\le C/R\). The small-set property applied
to \(A=\{V_i>R\}\) gives uniform integrability.

Under uniform integrability, choose \(R_j\ge j\) strictly increasing
such that
\[
 \sup_i\mathbb E[V_i\mathbf 1_{\{V_i>R_j\}}]\le 2^{-j}.
\]
With \(x^+=\max(x,0)\), put
\[
 \Phi(t)=t+\sum_{j\ge1}(t-R_j)^+ .
\]
The sum is locally finite, and \(\Phi\) is nonnegative,
nondecreasing and convex. For \(t\ge2R_J\), each of the first \(J\)
summands is at least \(t/2\), so
\(\Phi(t)/t\ge1+J/2\). Also, monotone convergence gives
\[
 \mathbb E\Phi(V_i)\le C+
     \sum_{j\ge1}\mathbb E[V_i\mathbf 1_{\{V_i>R_j\}}]\le C+1.
\]
Finally, under the convex-function condition, for any \(\eta>0\)
there is \(R\) such that \(t\le\eta\Phi(t)\) for \(t>R\).
The supremum of the tails is then at most
\(\eta\sup_i\mathbb E\Phi(V_i)\), which proves uniform
integrability as \(\eta\) tends to zero.
\end{proof}

\noindent
For the classical criteria and the scalar case, see
\cite[Chapter 13]{Williams},
\cite[Lemma 5.10]{Kallenberg}, and
\cite[Chapter 3, Section 5, Theorem 5.6]{Shorack}.

\section{Vitali convergence in mean}
\label{p207:W32}

\begin{theorem}[Vitali]
Let \(X_n\) be Bochner integrable in a Banach space and let \(X\) be
a Banach-valued random variable. If \(X_n\to X\) in probability and
\((X_n)\) is uniformly integrable, then \(X\) has a Bochner-integrable
representative and \(\mathbb E\lVert X_n-X\rVert\to0\).
Conversely, convergence in \(L^1\) implies convergence in probability
and uniform integrability of the sequence.
\end{theorem}

\begin{proof}
Choose a subsequence \(X_{n_j}\) with
\(\mathbb P(\lVert X_{n_j}-X\rVert>2^{-j})\le2^{-j}\).
The union of these exceptional events for \(j\ge J\) has probability
at most \(\sum_{j\ge J}2^{-j}\), so the subsequence converges to \(X\)
almost surely. Its chosen strongly measurable representatives are
Cauchy on a measurable set, since the Cauchy condition uses only
countably many inequalities between their norm differences.
Take the pointwise limit there and set it to zero elsewhere. This
is a strongly measurable representative of \(X\).
Indeed, all the simple approximants used for the subsequence have
values in a common separable closed subspace, as do the chosen
representatives and their limit. In this subspace the limit is
Borel measurable, since distances to closed sets are continuous.
Choose a countable dense sequence. For each \(m\), approximate the
map by a nearest point among the first \(m\) points, resolving ties
by the least index. Each approximant is measurable and finite-valued,
and the distances to the chosen points decrease to zero.
Fatou's lemma and the \(L^1\) bound give
\(\mathbb E\lVert X\rVert\le\liminf_j\mathbb E\lVert X_{n_j}\rVert
<\infty\).

The family enlarged by \(X\) still satisfies the small-set
condition in Section~\ref{p207:W31}. For \(X\) itself this follows
by splitting \(\lVert X\rVert\) at a large threshold and using
its integrability. For a fixed \(a>0\) let
\(B_n=\{\lVert X_n-X\rVert>a\}\). Then
\[
 \mathbb E\lVert X_n-X\rVert
 \le a+\mathbb E[\lVert X_n\rVert\mathbf 1_{B_n}]
       +\mathbb E[\lVert X\rVert\mathbf 1_{B_n}].
\]
The probabilities of \(B_n\) tend to zero, so the last two terms
tend to zero. Let \(a\) tend to zero.

For the converse, Markov's inequality gives convergence in probability.
Convergence in \(L^1\) also gives a common \(L^1\) bound. The estimate
\[
 \mathbb E[\lVert X_n\rVert\mathbf 1_A]
 \le \mathbb E[\lVert X\rVert\mathbf 1_A]
       +\mathbb E\lVert X_n-X\rVert
\]
controls the integral uniformly on small sets for large \(n\).
Each of the finitely many remaining norms has small integral on
small sets, so one can take the smallest of their thresholds.
Section~\ref{p207:W31} gives uniform integrability.
\end{proof}

\noindent
Compare \cite[Theorem 5.12]{Kallenberg} and the scalar versions in
\cite{P105,D11}.

\section{Almost-sure convergence from a first-moment bound}
\label{p207:W33}

\begin{theorem}
Let \((X_n,\mathcal F_n)\) be a real submartingale with
\(\sup_n\mathbb E|X_n|<\infty\). Then \(X_n\) converges almost surely
to an integrable, \(\mathcal F_\infty\)-measurable random variable
\(X_\infty\), where
\(\mathcal F_\infty=\sigma(\bigcup_n\mathcal F_n)\).
The moment hypothesis is equivalent to
\(\sup_n\mathbb E X_n^+<\infty\).
\end{theorem}

\begin{proof}
Since \(\mathbb EX_n\ge\mathbb EX_0\),
\(\mathbb E|X_n|=2\mathbb EX_n^+-\mathbb EX_n\)
gives the nontrivial direction of the equivalence.
For any rational \(a<b\), write \(U_n(a,b)\) for the number
of completed upcrossings through time \(n\) as in
Section~\ref{p207:W28}. The upcrossing inequality gives a finite
bound independent of \(n\),
\[
 (b-a)\mathbb E U_n(a,b)
 \le \sup_k\mathbb E X_k^++|a|+\mathbb E[(X_0-a)^-].
\]
The counts are nondecreasing in \(n\). Monotone convergence
shows that the total number is finite almost surely for this
interval, hence for all rational intervals simultaneously.
If \(\liminf_n X_n<\limsup_n X_n\) on a path, choose rational
\(a<b\) strictly between the two limits. The path then goes below
\(a\) and above \(b\) infinitely often and completes infinitely
many upcrossings. Thus the extended-real limit exists almost
surely. Fatou's lemma gives
\(\mathbb E\liminf_n|X_n|\le\sup_n\mathbb E|X_n|<\infty\),
which excludes infinite limits almost surely and gives
integrability. The limit superior is \(\mathcal F_\infty\)-measurable.
Use its value where finite and set it to zero elsewhere to obtain
\(X_\infty\).
\end{proof}

\noindent
The doubling example in Section~\ref{p207:W30} shows that the
hypothesis alone does not give \(L^1\) convergence.
For the upcrossing proof, see also
\cite[Sections 11.2--11.5]{Williams} and
\cite[Theorem 9.19]{Kallenberg}.

\section{Uniform integrability and terminal closure}
\label{p207:W34}

\begin{theorem}[Terminal closure]
Let \((M_n,\mathcal F_n)\) be a uniformly integrable real martingale.
There is an integrable \(\mathcal F_\infty\)-measurable variable
\(M_\infty\) with \(M_n\to M_\infty\) almost surely and in \(L^1\), and
\[
 M_n=\mathbb E[M_\infty\mid\mathcal F_n]
 \qquad\text{almost surely for every }n.
\]
Conversely, for any integrable \(Y\), the martingale
\((\mathbb E[Y\mid\mathcal F_n])_{n\ge0}\) is uniformly integrable.
A uniformly integrable real submartingale has the same two modes of
convergence and satisfies
\(X_n\le\mathbb E[X_\infty\mid\mathcal F_n]\) almost surely.
\end{theorem}

\begin{proof}
Uniform integrability implies \(L^1\) boundedness.
Section~\ref{p207:W33} supplies the almost-sure limit and its
\(\mathcal F_\infty\)-measurable representative. Almost-sure convergence
implies convergence in probability, so Section~\ref{p207:W32} gives
\(L^1\) convergence. For each fixed \(n\), the conditional-expectation
contraction gives
\[
 \big\lVert\mathbb E[M_m-M_\infty\mid\mathcal F_n]\big\rVert_1
 \le \lVert M_m-M_\infty\rVert_1\longrightarrow0.
\]
For \(m\ge n\), the conditional expectation of \(M_m\) equals \(M_n\).
Thus the limit satisfies the claimed identity. For a submartingale
\(X_n\le \mathbb E[X_m\mid\mathcal F_n]\) when \(m\ge n\).
The negative part of
\(\mathbb E[X_\infty\mid\mathcal F_n]-X_n\) has expectation at most
\(\lVert X_m-X_\infty\rVert_1\) by the same contraction, so it vanishes.

For the converse, the tower property gives the martingale identity.
For any sub-sigma-algebra \(\mathcal H\subseteq\mathcal F\), conditional
absolute value control and positivity give, for every measurable \(A\),
\[
 \mathbb E[|\mathbb E[Y\mid\mathcal H]|\mathbf 1_A]
 \le R\mathbb P(A)
    +\mathbb E[|Y|\mathbf 1_{\{|Y|>R\}}].
\]
In this estimate, split \(|Y|\) at \(R\), bound its low part by \(R\),
and bound the integral of the conditional high part over \(A\) by
its integral over the whole space. This estimate and the common
\(L^1\) bound \(\mathbb E|Y|\) give uniform integrability by
Section~\ref{p207:W31}, even when \(\mathcal H\) ranges over all
sub-sigma-algebras.
\end{proof}

\noindent
The conditional-expectation laws used here are proved in \cite{P105}.
Compare \cite[Theorem 9.22]{Kallenberg}.

\section{Higher moments and norm convergence}
\label{p207:W35}

\begin{theorem}
Let \(M\) be a real martingale and \(1<p<\infty\).
If \(C=\sup_n\mathbb E|M_n|^p<\infty\), then \(M_n\to M_\infty\)
almost surely and in \(L^p\). The terminal representation
\(M_n=\mathbb E[M_\infty\mid\mathcal F_n]\) still holds.
\end{theorem}

\begin{proof}
For \(R>0\),
\(\mathbb E[|M_n|\mathbf 1_{\{|M_n|>R\}}]\le R^{1-p}C\),
so the family is uniformly integrable and
Section~\ref{p207:W34} applies. The process \((|M_n|)\) is a
nonnegative submartingale by conditional absolute value control.
The finite \(L^p\) maximal inequality in Section~\ref{p207:W27}
gives
\[
 \mathbb E\left[\left(\max_{0\le k\le n}|M_k|\right)^p\right]
 \le \left(\frac{p}{p-1}\right)^p \mathbb E|M_n|^p.
\]
By monotone convergence, \(Z=\sup_n|M_n|\) has finite \(p\)-th
moment. In particular \(|M_\infty|\le Z\) almost surely and
\(|M_n-M_\infty|^p\le 2^p Z^p\) almost surely. Dominated
convergence gives the \(L^p\) limit.
\end{proof}

\noindent
Compare \cite[Chapter 14]{Williams} and
\cite[Corollary 9.23]{Kallenberg}.

\section{Lévy's upward theorem}
\label{p207:W36}

\begin{theorem}
For increasing sigma-algebras \(\mathcal F_n\) and integrable real \(Y\),
\[
 \mathbb E[Y\mid\mathcal F_n]\longrightarrow
 \mathbb E[Y\mid\mathcal F_\infty]
 \quad\text{almost surely and in }L^1.
\]
If \(Y\in L^p\) for \(1<p<\infty\), the convergence also holds in \(L^p\).
\end{theorem}

\begin{proof}
The conditional expectations are a uniformly integrable martingale
by Section~\ref{p207:W34}. Let \(Z\) be its almost-sure and \(L^1\)
limit there. If \(A\in\bigcup_n\mathcal F_n\), then for all large
\(m\), \(\int_A\mathbb E[Y\mid\mathcal F_m]\,d\mathbb P
=\int_A Y\,d\mathbb P\). Passing to the \(L^1\) limit gives
\(\int_A Z\,d\mathbb P=\int_A Y\,d\mathbb P\).
The union of the increasing sigma-algebras is a \(\pi\)-system
containing \(\Omega\). The events on which these integrals agree
form a \(\lambda\)-system, by finite integrability and countable
additivity. The \(\pi\)--\(\lambda\) theorem extends agreement to
\(\mathcal F_\infty\). Since \(Z\) is \(\mathcal F_\infty\)-measurable,
this identifies it as the stated conditional expectation.
If \(Y\in L^p\), conditional Jensen gives the common bound
\(\mathbb E|\mathbb E[Y\mid\mathcal F_n]|^p\le\mathbb E|Y|^p\),
and Section~\ref{p207:W35} applies.
\end{proof}

\noindent
Compare \cite[Sections 14.1--14.2]{Williams}.

\section{Lévy's downward theorem}
\label{p207:W37}

\begin{theorem}
For decreasing sigma-algebras \(\mathcal G_n\) put
\(\mathcal G_\infty=\bigcap_n\mathcal G_n\). If \(Y\) is real
and integrable, then
\[
 \mathbb E[Y\mid\mathcal G_n]\longrightarrow
 \mathbb E[Y\mid\mathcal G_\infty]
 \quad\text{almost surely and in }L^1.
\]
\end{theorem}

\begin{proof}
Write \(Z_n=\mathbb E[Y\mid\mathcal G_n]\).
For any \(N\), the reversed, sign-changed sequence
\((-Z_{N-k})_{0\le k\le N}\) is a martingale with respect to the
increasing sigma-algebras \((\mathcal G_{N-k})_{0\le k\le N}\),
by the tower property.

By Lemma~\ref{p207:greedy}, the upcrossing count on a finite path
is the maximum length of an ordered list of crossing pairs.
Reversing pair order, replacing each buy--sale pair \((s_i,t_i)\)
by \((N-t_i,N-s_i)\), and negating the path gives a bijection
with upcrossing lists for \([-b,-a]\) on the reversed, sign-changed
path. Thus the corresponding counts are equal.
The upcrossing inequality in Section~\ref{p207:W28} now gives
\[
 \begin{split}
 (b-a)\mathbb E U_N(a,b;Z)
 &\le \mathbb E[(-Z_0+b)^+]+\mathbb E[(-Z_N+b)^-]\\
 &\le 2\mathbb E|Y|+2|b|.
 \end{split}
\]
The last bound uses the \(L^1\) contraction of conditional
expectation. Monotone convergence and the rational-interval argument
of Section~\ref{p207:W33} imply almost-sure existence of an
extended-real limit. The same contraction and Fatou's lemma
make it finite and integrable. Since \(Z_n\) for \(n\ge m\) is
\(\mathcal G_m\)-measurable, \(\limsup_n Z_n\) is
\(\mathcal G_m\)-measurable for each \(m\). Replacing its
infinite values by zero therefore yields a
\(\mathcal G_\infty\)-measurable representative \(Z\).
The family \((Z_n)\) is uniformly integrable by the fixed-input
estimate of Section~\ref{p207:W34}, so \(Z_n\to Z\) in \(L^1\) by
Section~\ref{p207:W32}. For \(A\in\mathcal G_\infty\),
\(\int_A Z_n\,d\mathbb P=\int_A Y\,d\mathbb P\) for every \(n\).
Taking limits gives the defining conditional-expectation identity.
\end{proof}

\noindent
Compare \cite[Section 14.4]{Williams} and \cite[Section 4.7]{Durrett}.

\section{Orthogonal martingale increments}
\label{p207:W38}

\begin{theorem}
Let \(M\) be a real martingale with \(M_n\in L^2\) for every \(n\).
Write \(D_k=M_{k+1}-M_k\). The \(D_k\) are pairwise orthogonal
in \(L^2\) and each is orthogonal to \(M_0\). In particular
\[
 \mathbb E M_n^2=\mathbb E M_0^2+\sum_{k<n}\mathbb ED_k^2 .
\]
The martingale is \(L^2\)-bounded if and only if
\(\sum_{k\ge0}\mathbb ED_k^2<\infty\). In this case it converges
almost surely and in \(L^2\), and its limit satisfies
\[
 \mathbb E M_\infty^2=\mathbb E M_0^2+\sum_{k\ge0}\mathbb ED_k^2,
 \qquad
 \mathbb E(M_\infty-M_n)^2=\sum_{k\ge n}\mathbb ED_k^2 .
\]
\end{theorem}

\begin{proof}
For \(i<j\), \(D_i\) is \(\mathcal F_j\)-measurable and
\(\mathbb E[D_j\mid\mathcal F_j]=0\).
First multiply \(D_j\) by the bounded variable
\(D_i\mathbf 1_{\{|D_i|\le R\}}\) and take expectation, obtaining zero.
As \(R\to\infty\), Cauchy--Schwarz shows that these products converge
to \(D_i D_j\) in \(L^1\). Thus \(\mathbb E[D_iD_j]=0\).
The same argument with the \(\mathcal F_j\)-measurable variable
\(M_0\) gives \(\mathbb E[M_0D_j]=0\).
Expand \(M_n=M_0+\sum_{k<n}D_k\) and take expectation of its
square. This proves the identity and the boundedness equivalence.
In the bounded case use Section~\ref{p207:W35} with \(p=2\).
The first limit identity follows by continuity of the \(L^2\) norm,
and the second by applying this continuity to
\(M_m-M_n=\sum_{n\le k<m}D_k\) as \(m\to\infty\).
\end{proof}

\noindent
Compare \cite[Section 12.1]{Williams}.
\chapter{Optional Sampling and Stopped Sums}

\section{The conditional inequality for bounded stopping times}\label{p207:W39}

\begin{theorem}[Bounded optional sampling for submartingales]
Let \(X\) be a real integrable submartingale and let
\(\sigma\le\tau\le N\) be stopping times, where \(N\) is finite
and deterministic. Then
\[
  \mathbb E[X_\tau\mid\mathcal F_\sigma]\ge X_\sigma
       \quad\text{almost surely},
  \qquad
  \mathbb E X_\sigma\le \mathbb E X_\tau .
\]
\end{theorem}
\begin{proof}
By Section~\ref{p207:W02}, \(X_\sigma\) and \(X_\tau\) are
measurable for their respective stopped \(\sigma\)-algebras and
integrable, the latter since their absolute values are bounded
by \(\sum_{k=0}^N|X_k|\).
Pathwise telescoping gives
\[
 X_\tau-X_\sigma
     =\sum_{k=0}^{N-1}1_{\{\sigma\le k<\tau\}}(X_{k+1}-X_k).
\]
For \(A\in\mathcal F_\sigma\), the event
\(B_k=A\cap\{\sigma\le k\}\cap\{\tau>k\}\) belongs to
\(\mathcal F_k\). The integral of \(1_{B_k}(X_{k+1}-X_k)\)
is nonnegative by testing the one-step conditional expectation.
Therefore the integral of \(X_\tau-X_\sigma\) over \(A\)
is nonnegative. To see this is the conditional inequality,
write the difference of its two sides as the integrable
\(\mathcal F_\sigma\)-measurable variable \(D\). By the test-event
identity its integral on every such \(A\) is nonnegative.
Use \(A=\{D<0\}\) to conclude \(D\ge0\) almost surely.
Testing at \(\Omega\) yields the expected-value inequality.
\end{proof}

See \cite[Section 10.10 and Appendix A14]{Williams},
\cite[Theorem 9.12]{Kallenberg}, and \cite[Theorem 4.4.1]{Durrett}.

\section{The conditional identity for bounded stopping times}\label{p207:W40}

\begin{theorem}[Bounded optional sampling for martingales]
Under the same stopping hypotheses, a real integrable martingale
\(M\) satisfies
\[
 \mathbb E[M_\tau\mid\mathcal F_\sigma]=M_\sigma
       \quad\text{almost surely},
 \qquad
 \mathbb E M_\tau=\mathbb E M_\sigma .
\]
\end{theorem}
\begin{proof}
Apply the preceding theorem to both \(M\) and \(-M\). Both are
submartingales, and linearity of conditional expectation makes
the two inequalities opposite. Taking expectations gives the
second assertion.
\end{proof}

\section{Unbounded optional sampling for uniformly integrable martingales}\label{p207:W41}

\begin{theorem}[Sampling a closed martingale]
Let \(M\) be a uniformly integrable real martingale, and let
\(\sigma\le\tau\) be stopping times, possibly infinite. Use
the \(\mathcal F_\infty\)-measurable integrable terminal
representative \(M_\infty\) from Section~\ref{p207:W34} as the
stopped value on the infinite layer. Then \(M_\sigma,M_\tau\)
are integrable and
\[
 \mathbb E[M_\tau\mid\mathcal F_\sigma]=M_\sigma
        \quad\text{almost surely},
 \qquad
 \mathbb E M_\tau=\mathbb E M_\sigma=\mathbb E M_0 .
\]
\end{theorem}
\begin{proof}
Put \(Y=M_\infty\), so \(M_n=\mathbb E[Y\mid\mathcal F_n]\)
almost surely for every \(n\). For an arbitrary stopping time
\(\rho\), Section~\ref{p207:W02} proves \(M_\rho\) is
\(\mathcal F_\rho\)-measurable with the indicated convention.
Conditional positivity gives \(|M_n|\le\mathbb E[|Y|\mid\mathcal F_n]\).
Since \(\{\rho=n\}\in\mathcal F_n\),
\[
 \mathbb E[|M_n|1_{\{\rho=n\}}]
       \le \mathbb E[|Y|1_{\{\rho=n\}}].
\]
Sum over all finite layers and add the infinite layer, where
\(M_\rho=Y\). By nonnegative countable additivity of the integral,
\(\mathbb E|M_\rho|\le\mathbb E|Y|<\infty\).

For \(A\in\mathcal F_\rho\), each \(A\cap\{\rho=n\}\) is in
\(\mathcal F_n\), so the integrals of \(M_\rho=M_n\) and \(Y\)
over it agree. On \(A\cap\{\rho=\infty\}\) they agree by
definition. Summing the integrals is legitimate since both
variables are integrable. Thus
\(M_\rho=\mathbb E[Y\mid\mathcal F_\rho]\) almost surely.
Use \(\mathcal F_\sigma\subseteq\mathcal F_\tau\) and the tower
property to prove the conditional identity. Testing on
\(\Omega\) yields equality of expectations with \(\mathbb E Y\),
which equals \(\mathbb E M_0\) by the time-zero closure formula.
\end{proof}

For comparison, see \cite[Appendix A14]{Williams}
and \cite[Sections 4.4 and 4.8]{Durrett}.

\section{Integrable stopping times and controlled increments}\label{p207:W42}

\begin{theorem}
Let \(M\) be a real martingale and \(\tau\) a stopping time
with \(\mathbb E\tau<\infty\). Suppose there is a finite \(C\ge0\)
such that
\(\mathbb E[|M_{k+1}-M_k|\mid\mathcal F_k]\le C\)
almost surely for each \(k\ge0\).
Then \(M_\tau\) is integrable and
\(\mathbb E M_\tau=\mathbb E M_0\).
In particular, the conclusion holds if
\(|M_{k+1}-M_k|\le C\) almost surely for every \(k\).
\end{theorem}
\begin{proof}
An integrable stopping time is finite almost surely, so we may
assign \(M_\tau=0\) on the null infinite layer. Since
\(\{\tau>k\}\in\mathcal F_k\), the test-event identity bounds
\(\mathbb E[1_{\{\tau>k\}}|\Delta M_k|]\) by
\(C\mathbb P(\tau>k)\). Tonelli and the identity
\(\tau=\sum_{k\ge0}1_{\{\tau>k\}}\) give
\[
 \mathbb E\sum_{k\ge0}1_{\{\tau>k\}}|\Delta M_k|
     \le C\sum_{k\ge0}\mathbb P(\tau>k)
     =C\mathbb E\tau.
\]
For every \(n\), \(M_{\tau\wedge n}\) is bounded in absolute
value by \(|M_0|\) plus this absolute sum. It converges almost
surely to \(M_\tau\), which is integrable by the same bound.
Dominated convergence transfers
\(\mathbb E M_{\tau\wedge n}=\mathbb E M_0\) from bounded optional
sampling. The bounded-increment special case follows from
conditional positivity.
\end{proof}

\begin{example}[Why the hypotheses matter]
For simple symmetric walk on \(\mathbb Z\) starting at zero,
the first hitting time of \(+1\) is finite almost surely and has
infinite expectation by the symmetric one-dimensional case of
Section~\ref{p207:W21}, after reflection. The terminal walk value
is \(1\) almost surely and its initial value is zero.
Thus an almost surely finite stop and bounded increments alone
can fail to preserve the expectation.

For the fair-coin doubling martingale from
Section~\ref{p207:W30}, let \(\tau\) be the first tail, numbered
starting at one. It is a stopping time with
\(\mathbb P(\tau>k)=2^{-k}\) for \(k\ge0\), so \(\mathbb E\tau=2\).
Its stopped value is zero almost surely, whereas \(M_0=1\).
On the event of \(k\) consecutive heads, the conditional expected
absolute increment at time \(k\) is \(2^k\); this event has
positive probability. This time no uniform bound on the
conditional expected absolute increment is available, and an
integrable stop alone fails.
\end{example}

\section{Wald's first and second identities}\label{p207:W43}

\begin{theorem}
Let \(Y_1,Y_2,\ldots\) be identically distributed integrable real
variables, with \(Y_k\) \(\mathcal F_k\)-measurable and independent of
\(\mathcal F_{k-1}\). Put \(\mu=\mathbb E Y_1\),
\(S_n=\sum_{k=1}^nY_k\), \(S_0=0\). If \(\tau\) is a stopping time
with \(\mathbb E\tau<\infty\), then \(S_\tau\) is integrable and
\[
 \mathbb E S_\tau=\mu\,\mathbb E\tau.
\]
If in addition \(v=\mathbb E[(Y_1-\mu)^2]<\infty\), then
\[
  \mathbb E[(S_\tau-\mu\tau)^2]=v\,\mathbb E\tau.
\]
\end{theorem}
\begin{proof}
By independence \(M_n=S_n-\mu n\) is a martingale, and its
conditional expected absolute increment is
\(\mathbb E|Y_1-\mu|\) at every time. The preceding theorem
gives \(M_\tau\in L^1\) and \(\mathbb E M_\tau=0\).
Since \(\mu\tau\) is integrable, this proves the first identity.

For the second, put \(Z_k=1_{\{\tau\ge k\}}(Y_k-\mu)\) for \(k\ge1\).
The event is \(\mathcal F_{k-1}\)-measurable and independence gives
\(\mathbb E[Z_k\mid\mathcal F_{k-1}]=0\) and
\(\mathbb E Z_k^2=v\mathbb P(\tau\ge k)\).
Thus the finite sums \(\sum_{k=1}^n Z_k=M_{n\wedge\tau}\) form a
square-integrable martingale with orthogonal increments, as in
Section~\ref{p207:W38}. The sum of their increment second moments is
\(v\sum_{k\ge1}\mathbb P(\tau\ge k)=v\mathbb E\tau\).
Orthogonality makes the partial sums Cauchy in \(L^2\).
They also converge almost surely to \(M_\tau\) because \(\tau\)
is finite almost surely. The \(L^2\) limit agrees almost surely
with this pathwise limit, by uniqueness of limits in probability.
Take the limit in
\(\mathbb E M_{n\wedge\tau}^2=v\sum_{k=1}^n\mathbb P(\tau\ge k)\)
to obtain the claimed second moment.
\end{proof}

\backmatter
\begingroup
\small
\begin{thebibliography}{99}
\setlength{\itemsep}{0.4em}
\bibitem{Norris} J. R. Norris, \emph{Markov Chains},
  Cambridge University Press, 1997.
\bibitem{Williams} D. Williams, \emph{Probability with Martingales},
  Cambridge University Press, 1991.
\bibitem{Durrett} R. Durrett, \emph{Probability: Theory and Examples},
  5th ed., Cambridge University Press, 2019.
\bibitem{Kallenberg} O. Kallenberg, \emph{Foundations of Modern Probability},
  3rd ed., Springer, 2021.
\bibitem{Knill} O. Knill, \emph{Probability Theory and Stochastic Processes
  with Applications}, Section 3.8,
  \href{https://people.math.harvard.edu/~knill/books/KnillProbability.pdf}
       {author's copy}.
\bibitem{Shorack} G. R. Shorack, \emph{Probability for Statisticians},
  Chapter 3, Section 5, Theorem 5.6,
  \href{https://sites.stat.washington.edu/jaw/COURSES/520s/521/bk521reJaw2012.pdf}
       {University of Washington course edition, 2012}.
\bibitem{Azuma} K. Azuma, Weighted sums of certain dependent random
  variables, \emph{Tohoku Math. J.} 19 (1967), 357--367.
  \href{https://doi.org/10.2748/tmj/1178243286}{DOI}.
\bibitem{D05} MathGenome, \emph{Linear Algebra and Foundational Higher
  Algebra}, D05 unified notes, Module 13, real spectral theorem.
\bibitem{D11} MathGenome, \emph{Measure, Integration and Real Analysis},
  D11 unified notes, measure extension, integration, and Vitali convergence.
\bibitem{P105} MathGenome, \emph{Probability: Common Foundations and Limit
  Theorems}, P1-05 unified notes,
  uniform integrability, conditional expectation, path measures,
  and laws of large numbers.
\bibitem{P205} MathGenome, \emph{Extremal Methods, Networks, and
  Matroid Optimization}, P2-05 unified notes, W25.
\bibitem{Mathlib} The Mathlib Community, \emph{Mathlib},
  \href{https://github.com/leanprover-community/mathlib4/tree/56a005855ea17e0b0958f73ec801ae140b0a1679}
       {fixed revision \texttt{56a005855ea1}}.
\end{thebibliography}
\endgroup
\end{document}
